% !TeX root = ../../Book3.tex
% 纲要标题：## 第25章*　Einstein 方程、规范与局部模问题
% 纲要位置：几何学计划纲要.md，第 3165 行。
% 正文状态：初稿；尚未正式审核。

% 纲要内容开始
% 
% > 本章目标与证明范围：与 Yamabe 分章：Einstein 方程是模微分同胚的系统，而非标量共形方程。完整建立线性化、规范固定和有限维局部障碍模型；一般存在与分类不在本章无条件承诺。 非线性共同输入定位于 III.13.5—13.6，不以完整阅读 Yamabe 章为先修。
% 
% 纲要内容结束

\OptionalChapter{Einstein 方程、规范与局部模问题}{g3:ch:25}
\BookChapterNumber{g3:ch:25}

\begin{chapterabstract}
	Einstein 方程要求 Ricci 张量与度量成比例，而微分同胚使同一个几何对象具有无穷多个张量代表。本章从度量空间的切向量出发，推导联络、Ricci 张量、标量曲率及总体积的变分公式。随后用守恒的 Einstein 张量和背景散度构造椭圆方程，证明规范项在附近的解上消失，并通过核正交椭圆方程建立局部切片。形式自伴的线性化给出有限维无穷小形变空间；隐函数定理消去其正交补，剩下的有限维障碍映射及等距群作用描述局部模问题。空间形式、乘积和平坦环面分别说明刚性、可积形变、商空间与整体退化，Kähler 类中的 Ricci 平坦存在定理则作为明确的外部结果介绍。
\end{chapterabstract}

\begin{learninggoals}
	\begin{itemize}
		\item 在本卷曲率与 Laplace 符号约定下推导曲率线性化，并识别微分同胚产生的核方向。
		\item 计算带体积约束的 Einstein--Hilbert 泛函的一阶变分，区分固定常数、固定体积和模缩放问题。
		\item 验证规范化方程的椭圆性，证明规范等价与散度切片，并说明 Killing 场留下的剩余对称。
		\item 用 Fredholm 理论与隐函数定理构造有限维障碍映射，计算二阶障碍并判断局部刚性。
		\item 在空间形式、乘积和平坦环面上检验一般结论，区分局部参数、群作用和整体紧性。
	\end{itemize}
\end{learninggoals}

设 \(g_0\) 是圆球上的标准度量，\(\phi_t\) 是一族微分同胚。张量 \(\phi_t^*g_0\) 的坐标系数通常随 \(t\) 改变，但各个度量都与 \(g_0\) 等距。这种变化在微分后成为 \(\mathcal L_Xg_0\)。若把它当作真正改变几何的未知方向，曲率方程的线性化便出现一个由任意向量场产生的无限维核。研究附近的几何结构，首先要把这种重参数化自由处理清楚。

\ProseParagraphBreak
本章研究的条件是 \(\operatorname{Ric}_g=\lambda g\)，称为 Einstein 方程。\footnote{Albert Einstein（1879--1955），德国出生的理论物理学家，建立了广义相对论。}在维数至少为三的连通流形上，缩并 Bianchi 恒等式迫使这里的比例因子为常数。固定共形类的标量曲率问题只改变一个正函数；这里允许全部对称二张量方向变化，方程是一个带重参数化对称的系统。前一章的临界极小化论证不是本章的先修；直接使用的非线性工具来自第13.5—13.6节。

\ProseParagraphBreak
我们从一个已知的闭 Einstein 流形 \((M^n,g_0)\) 出发，问附近的解能否用有限个参数描述。先固定 Einstein 常数 \(\lambda\)，用散度条件选择代表，再用守恒恒等式证明所加的规范项确实消失。椭圆线性化的核与余核有限维，因而可以把无限维方程分解为一个可逆方向上的方程和一个有限维方程。后者可能有非零高阶项，所以无穷小形变不一定能延伸成真正的度量族；即使障碍消失，还须除去背景度量的等距群。

\ProseParagraphBreak
阅读时需要熟悉第1—3章的联络与曲率计算，以及第11—13章的 Sobolev、椭圆和 Banach 空间理论。正 Einstein 常数下的迹消去会调用第16章的第一特征值估计。固定常数、整体缩放与体积归一化的关系将在第25.3节末明确处理；最后一节用模型检验局部结论，并说明一般存在性与整体紧性还需要什么几何信息。

% !TeX root = ../../Book3.tex
% 纲要标题：### 25.1 度量空间和曲率变分
% 纲要位置：几何学计划纲要.md，第 3169 行。
% 正文状态：初稿；尚未正式审核。

% 纲要内容开始
% 
% * 对称二张量的切空间
% * Ricci 与标量曲率的线性化
% * 微分同胚方向的退化
% * Einstein–Hilbert 泛函和体积约束
% 
% 纲要内容结束

\section{度量空间和曲率变分}
\label{g3:sec:2501}

把度量的各个分量都作为未知函数，Ricci 曲率便成为一个含二阶导数的非线性算子。它和标量曲率方程有一个重要差别：改变坐标会同时改变全部度量分量，却不改变几何。为了看清这种退化，本节先在固定流形上求度量、联络和曲率的一阶变化，再把坐标变化放进同一公式。积分变分随后给出一个具有散度恒等式的 Einstein 方程；这一恒等式将在下一节帮助检验规范固定后的解。

以下积分均使用黎曼体积密度，因而不要求流形定向。除局部计算外，始终考虑闭连通光滑流形 \(M^n\)，\(n\ge3\)。对一个指定的光滑度量 \(g\)，\(\nabla\)、\(\operatorname{Ric}_g\)、\(s_g\) 分别表示其 Levi--Civita 联络、Ricci 张量和标量曲率。

\subsection{对称二张量的切空间}
\label{g3:subsec:2501:01}
在每一点，正定对称矩阵组成对称矩阵空间中的开凸锥。因此度量的变化方向不必正定；正定性限制的是变化后的度量，而不是切向量。

\begin{definition}\label{g3:def:einstein-metric-space}
	设 \(M^n\) 为闭连通光滑流形，\(n\ge3\)，\(g_0\) 为一个光滑黎曼度量，\(s>n/2+2\) 为整数。记
	\[
		\mathcal M^s(M)=\{g\in H^s(M;\operatorname{Sym}^2\mathrm T^*M):g(x)>0\text{ 对每个 }x\in M\}.
	\]
	其中 Sobolev 范数用 \(g_0\) 及其联络定义；连续代表由 Sobolev 嵌入唯一确定。若一个光滑度量 \(g\) 满足
	\begin{equation}\label{g3:eq:einstein-equation-fixed-lambda}
		\operatorname{Ric}_g=\lambda g
	\end{equation}
	且 \(\lambda\) 为实常数，就称 \(g\) 为\term[Einstein duliang]{Einstein 度量}[Einstein metric]，称 \(\lambda\) 为其 Einstein 常数。
\end{definition}

\begin{proposition}\label{g3:prop:metric-sobolev-open}
	设 \(M^n\) 为闭连通光滑流形，\(n\ge3\)，\(g_0\) 为光滑黎曼度量，\(s>n/2+2\) 为整数。则 \(\mathcal M^s(M)\) 是 \(H^s(M;\operatorname{Sym}^2\mathrm T^*M)\) 的开集，在任意 \(g\in\mathcal M^s(M)\) 处的切空间都是该对称二张量空间。映射
	\[
		g\longmapsto g^{-1},\qquad g\longmapsto\mu_g/\mu_{g_0}
	\]
	分别是到相应 \(H^s\) 截面空间的光滑映射。Ricci 张量与标量曲率给出光滑映射
	\[
		\operatorname{Ric}:\mathcal M^s(M)\longrightarrow H^{s-2}(M;\operatorname{Sym}^2\mathrm T^*M),
		\qquad
		\operatorname{Scal}:\mathcal M^s(M)\longrightarrow H^{s-2}(M).
	\]
\end{proposition}
\begin{proof}
	紧性使连续正定 \(g\) 相对于 \(g_0\) 的最小特征值有正下界 \(c\)。若 \(\|h\|_{C^0,g_0}<c\)，则 \(g+h\) 仍正定。第十一章的 Sobolev 嵌入给出 \(\|h\|_{C^0}\le C\|h\|_{H^s}\)，故正定条件在 \(H^s\) 中开。对每个对称张量 \(h\)，直线 \(g+th\) 在充分小的 \(|t|\) 上属于此开集，因而没有额外的切向限制。

	正定矩阵的求逆与行列式平方根是光滑纤维映射；在 \(g\) 的一个小邻域内，特征值一致远离零。将\cref{g3:thm:sobolev-composition-smooth}用于这些纤维映射，便得前两个光滑性结论。联络系数由 \(g^{-1}\partial g\) 的线性组合给出，曲率由 \(g^{-1}\partial^2g\) 及两个一阶导数的乘积给出，都是二阶喷射的光滑函数。由于 \(s-2>n/2\)，\cref{g3:prop:nonlinear-jet-operator}正好给出到 \(H^{s-2}\) 的光滑性。各图上的公式来自同一个张量，故按有限坐标覆盖粘合。
\end{proof}

这个开集描述也说明，光滑计算可以直接给出 Sobolev 映射的 Fréchet 导数：先在光滑方向上求导，再用已证的光滑性与稠密性延拓。所需的是固定源上的张量运算；变化坐标产生的复合另有导数计数。

\begin{lemma}\label{g3:lem:metric-inverse-volume-variation}
	设 \(M^n\) 为闭光滑流形，\(g(t)\) 为一条光滑黎曼度量曲线，\(g(0)=g\)、\(\dot g(0)=h\)，并记 \(H=\operatorname{tr}_g h\)。则
	\begin{equation}\label{g3:eq:inverse-volume-variation}
		\left.\frac{\mathrm d}{\mathrm dt}\right|_0g(t)^{ij}=-h^{ij},
		\qquad
		\left.\frac{\mathrm d}{\mathrm dt}\right|_0\mu_{g(t)}=\frac12H\mu_g.
	\end{equation}
	因此 \(V(g)=\int_M\mu_g\) 满足
	\[
		\mathrm DV_g(h)=\frac12\int_MH\,\mu_g.
	\]
\end{lemma}
\begin{proof}
	微分 \(g(t)^{ia}g(t)_{aj}=\delta^i_j\)，得到逆矩阵公式。体积密度在坐标中是 \(\sqrt{\det g(t)}\,|\mathrm dx^1\cdots\mathrm dx^n|\)。行列式的导数公式给出
	\[
		\left.\frac{\mathrm d}{\mathrm dt}\right|_0\log\det g(t)=\operatorname{tr}(g^{-1}h)=H,
	\]
	再微分平方根便得第二式。紧性允许在积分号下求导。
\end{proof}

例如 \(g(t)=(1+t)g\) 的切向量是 \(h=g\)，体积变化率为 \(nV(g)/2\)。固定体积的切向条件则只要求 \(\int H\,\mu_g=0\)，并不要求 \(H\) 在每一点为零；这两个条件将在形变计算中承担不同职责。

\subsection{Ricci 与标量曲率的线性化}
\label{g3:subsec:2501:02}
联络变化含一个度量导数，曲率变化再多一个导数。将这两步分开，可以保留协变形式，也能避免把指标升降中的变化混入曲率本身。

\begin{lemma}\label{g3:lem:levi-civita-variation}
	设 \(g(t)\) 为光滑流形 \(M\) 上的一条光滑黎曼度量曲线，\(g(0)=g\)、\(\dot g(0)=h\)。令
	\[
		C(X,Y)=\left.\frac{\mathrm d}{\mathrm dt}\right|_0\nabla^{g(t)}_XY
	\]
	其中 \(X,Y\) 固定而不随 \(t\) 变化。则 \(C\) 是对称的 \((1,2)\) 型张量，且
	\begin{equation}\label{g3:eq:connection-variation}
		2g(C(X,Y),Z)=(\nabla_Xh)(Y,Z)+(\nabla_Yh)(X,Z)-(\nabla_Zh)(X,Y).
	\end{equation}
	在坐标中，
	\[
		C^k_{ij}=\frac12g^{ka}\bigl(\nabla_i h_{ja}+\nabla_j h_{ia}-\nabla_a h_{ij}\bigr).
	\]
\end{lemma}
\begin{proof}
	两条联络之差是张量，微分仍保留张量性；无挠性给出 \(C(X,Y)=C(Y,X)\)。微分度量相容式，整理出
	\[
		(\nabla_Xh)(Y,Z)=g(C(X,Y),Z)+g(Y,C(X,Z)).
	\]
	将 \((X,Y,Z)\) 分别换成 \((Y,X,Z)\)、\((Z,X,Y)\)，把前两式相加再减去第三式。利用 \(C\) 的对称性，其余项相消，得到所需公式。
\end{proof}

下面固定所有收缩和伴随的符号。对对称二张量 \(h\) 与一形式 \(\omega\)，定义
\begin{equation}\label{g3:eq:divergence-bianchi-conventions}
	(\delta_g h)_j=-\nabla^i h_{ij},\qquad
	(\delta_g^*\omega)_{ij}=\frac12(\nabla_i\omega_j+\nabla_j\omega_i),\qquad
	\beta_g h=\delta_g h+\frac12\mathrm d\operatorname{tr}_g h.
\end{equation}
在一形式上，\(\delta_g\omega=-\nabla^i\omega_i\)，所以 \(\delta_g\mathrm df=\Delta_g f\)。闭流形上的分部积分说明 \(\delta_g^*\) 正是 \(\delta_g\) 的 \(L^2(g)\) 伴随。\(\beta_g\) 称为 \term[Bianchi suanzi]{Bianchi 算子}[Bianchi operator]；本节先用它组织曲率的导数。

对一点处的正交单位标架 \((e_a)\)，记
\begin{equation}\label{g3:eq:curvature-action-symmetric-tensors}
	\begin{aligned}
		\bigl(\mathring R_g(h)\bigr)_{ij}
		&=\sum_{a,b}h_{ab}\,g(R(e_a,e_i)e_j,e_b),\\
		\Delta_Lh&=\nabla^*\nabla h+\operatorname{Ric}_g\circ h+h\circ\operatorname{Ric}_g-2\mathring R_g(h).
	\end{aligned}
\end{equation}
这里 \((\operatorname{Ric}_g\circ h)_{ij}=\operatorname{Ric}_i{}^a h_{aj}\)，\(\nabla^*\nabla\) 使用非负主符号；后一算子称为\term[Lichnerowicz Laplace suanzi]{Lichnerowicz Laplace 算子}[Lichnerowicz Laplacian]。曲率的对称性保证这些定义与正交标架无关，并给出
\[
	\mathring R_g(g)=\operatorname{Ric}_g,\qquad
	\operatorname{tr}_g\mathring R_g(h)=\langle\operatorname{Ric}_g,h\rangle_g,
	\qquad \operatorname{tr}_g\Delta_Lh=\Delta_g\operatorname{tr}_g h.
\]
特别地，\(\Delta_Lg=0\)。这些检验固定了曲率作用中的槽位顺序。

\begin{theorem}\label{g3:thm:ricci-scalar-linearization}
	设 \(g(t)\) 为光滑流形 \(M\) 上的一条光滑黎曼度量曲线，\(g(0)=g\)、\(\dot g(0)=h\)，\(H=\operatorname{tr}_g h\)。曲率符号采用本书约定，\(\delta_g,\beta_g,\Delta_L\) 由式\eqref{g3:eq:divergence-bianchi-conventions}、\eqref{g3:eq:curvature-action-symmetric-tensors}定义。则
	\begin{align}
		\mathrm D\operatorname{Ric}_g(h)&=\frac12\Delta_Lh-\delta_g^*\beta_g h,\label{g3:eq:ricci-linearization}\\
		\mathrm D\operatorname{Scal}_g(h)&=\Delta_gH+\delta_g\delta_g h-\langle\operatorname{Ric}_g,h\rangle_g.\label{g3:eq:scalar-linearization}
	\end{align}
\end{theorem}
\begin{proof}
	先保持 Ricci 张量的两个指标都在下方。微分曲率坐标式\eqref{g3:eq:curvature-coordinate-components}，或在一点取 \(g\) 的法坐标，得
	\[
		(\mathrm D\operatorname{Ric}_g h)_{ij}=\nabla_k C^k_{ij}-\nabla_i C^k_{kj}.
	\]
	由 \(C^k_{kj}=\tfrac12\nabla_jH\)，得到尚未交换协变导数的表达
	\begin{equation}\label{g3:eq:ricci-linearization-uncommuted}
		2(\mathrm D\operatorname{Ric}_g h)_{ij}
		=\nabla_k\nabla_i h_j{}^k+\nabla_k\nabla_j h_i{}^k-\nabla_k\nabla^k h_{ij}-\nabla_i\nabla_jH.
	\end{equation}
	这里 \(h_j{}^k=g^{ka}h_{ja}\)，升指标使用固定的 \(g\)。对一个上指标和一个下指标，协变导数的交换分别给出正、负曲率作用，因而
	\[
		\nabla_k\nabla_i h_j{}^k
		=\nabla_i\nabla_k h_j{}^k+\operatorname{Ric}_i{}^a h_{aj}-\bigl(\mathring R_g(h)\bigr)_{ij}.
	\]
	另一项交换 \(i,j\) 后相同。将 \(\nabla_kh_j{}^k=-(\delta_gh)_j\) 代回式\eqref{g3:eq:ricci-linearization-uncommuted}，便有
	\[
		2\mathrm D\operatorname{Ric}_g(h)
		=\Delta_Lh-2\delta_g^*\delta_gh-\operatorname{Hess}_gH
		=\Delta_Lh-2\delta_g^*\beta_gh.
	\]
	因为 \(\delta_g^*\mathrm dH=\operatorname{Hess}_gH\)，得到第一式。

	标量曲率是 \(g^{ij}\operatorname{Ric}_{ij}\)，所以还要微分逆度量。由式\eqref{g3:eq:inverse-volume-variation}，
	\[
		\mathrm D\operatorname{Scal}_g(h)
		=-\langle\operatorname{Ric}_g,h\rangle_g+\operatorname{tr}_g\mathrm D\operatorname{Ric}_g(h).
	\]
	前面的迹恒等式以及 \(\operatorname{tr}_g\delta_g^*\omega=-\delta_g\omega\) 给出
	\[
		\operatorname{tr}_g\mathrm D\operatorname{Ric}_g(h)
		=\tfrac12\Delta_gH+\delta_g\bigl(\delta_gh+\tfrac12\mathrm dH\bigr)
		=\Delta_gH+\delta_g\delta_gh.
	\]
	这就是第二式。
\end{proof}

\begin{example}\label{g3:exm:ricci-conformal-linearization}
	设 \(g\) 为光滑黎曼度量，\(f\) 为光滑实函数，取 \(h=fg\)。此时
	\[
		\delta_gh=-\mathrm df,\qquad \beta_gh=\frac{n-2}{2}\mathrm df,\qquad \Delta_L(fg)=(\Delta_gf)g.
	\]
	因此
	\begin{align*}
		\mathrm D\operatorname{Ric}_g(fg)&=\frac12(\Delta_gf)g-\frac{n-2}{2}\operatorname{Hess}_gf,\\
		\mathrm D\operatorname{Scal}_g(fg)&=(n-1)\Delta_gf-s_gf.
	\end{align*}
	这与第3.4节共形曲率公式在 \(e^{tf}g\) 上的导数一致。若 \(f=c\) 为常数，则 Ricci 的导数为零，而标量曲率的导数为 \(-cs_g\)：常数缩放不改变作为协变二张量的 Ricci，却会改变取迹所用的逆度量。
\end{example}

\subsection{微分同胚方向的退化}
\label{g3:subsec:2501:03}
对一个自微分同胚 \(\phi:M\to M\)，度量 \(g\) 与 \(\phi^*g\) 描述同一个几何结构，只是点的标记发生了变化。对曲线求导后，这种对称性成为 Einstein 线性化的核；在最高阶项中，它也表现为主符号的退化。

\begin{proposition}\label{g3:prop:einstein-diffeomorphism-degeneracy}
	设 \((M^n,g_0)\) 为闭光滑 Einstein 流形，\(n\ge3\)，\(\operatorname{Ric}_{g_0}=\lambda g_0\)。定义
	\(\mathcal E_\lambda(g)=\operatorname{Ric}_g-\lambda g\)。若 \(X\) 为光滑向量场，\(\omega=X^\flat\)，则
	\begin{equation}\label{g3:eq:einstein-gauge-linear-kernel}
		\mathcal L_Xg_0=2\delta_{g_0}^*\omega,\qquad
		\mathrm D\mathcal E_\lambda|_{g_0}(\mathcal L_Xg_0)=0.
	\end{equation}
	对任意非零实余向量 \(\xi\)，采用 \(\partial_j\mapsto i\xi_j\) 的主符号约定，有
	\begin{equation}\label{g3:eq:einstein-ungauged-symbol}
		\begin{aligned}
			\sigma_2(\mathrm D\mathcal E_\lambda)(\xi)h
			=\frac12\bigl(&|\xi|^2h+\xi\otimes\xi\operatorname{tr}h\\
			&-\xi\otimes h(\xi^\sharp,\cdot)-h(\xi^\sharp,\cdot)\otimes\xi\bigr).
		\end{aligned}
	\end{equation}
	其核恰为 \(\{\xi\otimes v+v\otimes\xi:v\in\mathrm T^*_xM\}\)，故未固定规范的 Einstein 方程不是椭圆系统。
\end{proposition}
\begin{proof}
	由 Lie 导数定义及联络的度量相容性，
	\[
		(\mathcal L_Xg_0)(Y,Z)=g_0(\nabla_YX,Z)+g_0(Y,\nabla_ZX),
	\]
	得到第一式。设 \(\phi_t\) 是 \(X\) 的局部流；闭性保证可以共同选取一个时间区间。拉回联络仍满足无挠与度量相容条件，Levi--Civita 联络的唯一性因此给出曲率的自然性，从而
	\[
		\mathcal E_\lambda(\phi_t^*g_0)=\phi_t^*\mathcal E_\lambda(g_0)=0.
	\]
	在 \(t=0\) 微分，便得第二式。

	在一点的法坐标中，只保留式\eqref{g3:eq:ricci-linearization-uncommuted}的二阶部分，并把 \(\partial_j\) 换成 \(i\xi_j\)，即得式\eqref{g3:eq:einstein-ungauged-symbol}。令 \(|\xi|=1\)、\(\xi=e^1\)。输出的 \((a,b)\) 分量在 \(a,b\ge2\) 时等于 \(h_{ab}/2\)，故核中的张量必须满足所有这些分量为零。反过来，任意仅有 \(h_{11}\)、\(h_{1a}\) 的对称张量都可写成 \(\xi\otimes v+v\otimes\xi\)，其中 \(v_1=h_{11}/2\)、\(v_a=h_{1a}\)；代入符号公式，其各项相消。一般 \(\xi\) 由齐次性恢复。
\end{proof}

主符号的核并非来自某个特殊曲率，而是来自任意坐标变化。即使某个 Einstein 度量没有非零 Killing 场，任意向量场仍给出式\eqref{g3:eq:einstein-gauge-linear-kernel}中的核方向；Killing 场只是其中使 \(\mathcal L_Xg_0=0\) 的向量场。下一节将分别处理需要去掉的轨道方向和作用本身的稳定子。

在函数空间中还须分清两个光滑性问题。对固定的光滑 \(\phi\)，\(g\mapsto\phi^*g\) 是 \(H^s\) 上的有界线性同构；若 \(\phi\) 也变化，导数 \(\mathcal L_Xg\) 会出现 \(X^a\partial_ag_{ij}\)，一般只由 \(g\in H^s\) 得到 \(H^{s-1}\) 控制。因而不能把整个微分同胚作用直接当作同一个 \(H^s\) 空间上的光滑作用。规范方程必须按所用源空间、目标空间和导数阶数另行建立。

\subsection{Einstein–Hilbert 泛函和体积约束}
\label{g3:subsec:2501:04}
方程 \(\operatorname{Ric}=\lambda g\) 不只来自局部曲率条件，也来自总标量曲率的驻值。与仅改变共形因子的变分不同，现在允许全部对称二张量方向，因而 Euler 方程保留整个 Ricci 张量。

定义\term[Einstein Hilbert fanhan]{Einstein--Hilbert 泛函}[Einstein--Hilbert functional]及其体积归一化形式
\begin{equation}\label{g3:eq:einstein-hilbert-functionals}
	\mathcal H(g)=\int_Ms_g\,\mu_g,\qquad
	\mathcal S(g)=V(g)^{-(n-2)/n}\mathcal H(g).
\end{equation}
后一个泛函在常数缩放下不变，因为 \(V(cg)=c^{n/2}V(g)\)、\(s_{cg}=c^{-1}s_g\)。这与第24章在共形类上使用的几何泛函相同，但这里的竞争度量不受共形限制。

\begin{theorem}\label{g3:thm:einstein-hilbert-critical}
	设 \(M^n\) 为闭连通光滑流形，\(n\ge3\)，\(g\) 为光滑黎曼度量。对任意光滑对称二张量 \(h\)，有
	\begin{equation}\label{g3:eq:einstein-hilbert-first-variation}
		\mathrm D\mathcal H_g(h)=-\int_M\left\langle\operatorname{Ric}_g-\frac12s_gg,h\right\rangle_g\,\mu_g.
	\end{equation}
	以下三项等价：\(g\) 是 Einstein 度量；\(g\) 是 \(\mathcal H\) 在体积 \(V(g)\) 的度量中关于光滑变分的临界点；\(g\) 是 \(\mathcal S\) 关于全部光滑变分的临界点。
\end{theorem}
\begin{proof}
	将式\eqref{g3:eq:scalar-linearization}与体积变化公式相加，得到
	\[
		\mathrm D\mathcal H_g(h)=\int_M\left(\Delta_gH+\delta_g\delta_gh-\langle\operatorname{Ric}_g,h\rangle_g+\frac12s_gH\right)\mu_g.
	\]
	前两项都是散度，闭流形上积分为零，故得到第一变分公式。

	体积约束的切向量满足 \(\int_MH\,\mu_g=0\)。每个这样的光滑 \(h\) 都由一条精确保体积的曲线实现：令
	\[
		\widetilde g(t)=\left(\frac{V(g)}{V(g+th)}\right)^{2/n}(g+th).
	\]
	在充分小的时间区间内它正定，体积恒为 \(V(g)\)，且因为 \(\mathrm DV_g(h)=0\)，有 \(\dot{\widetilde g}(0)=h\)。若 \(g\) 为约束临界点，\(\mathrm D\mathcal H_g\) 在 \(\ker\mathrm DV_g\) 上为零，而 \(\mathrm DV_g(g)=nV(g)/2\ne0\)。将任意 \(h\) 减去适当的常数倍 \(g\)，可知存在一个实常数 \(c\)，使 \(\mathrm D\mathcal H_g=c\,\mathrm DV_g\)。第一变分公式于是给出
	\[
		\operatorname{Ric}_g-\frac12s_gg=-\frac c2g.
	\]
	两边取迹，得 \((n-2)s_g=nc\)，所以 \(s_g\) 为常数，且 \(\operatorname{Ric}_g=(s_g/n)g\)。反过来，若 \(\operatorname{Ric}_g=\lambda g\)，则第一变分是 \((n-2)\lambda\mathrm DV_g\)，在体积约束切空间上为零。

	最后，\(\mathcal S\) 在缩放方向的导数恒为零；在体积约束切方向上，它的导数等于正数 \(V(g)^{-(n-2)/n}\) 乘以 \(\mathrm D\mathcal H_g\)。每个方向都唯一分成一个体积约束切方向与一个常数倍 \(g\)，故两种临界点条件等价。
\end{proof}

无体积约束的 \(\mathcal H\) 临界点则满足 \(\operatorname{Ric}-s g/2=0\)。在 \(n\ge3\) 时取迹即得 \(s=0\)，所以它们恰是 Ricci 平坦度量。二维的迹方程失去这一作用；本章的维数条件正是区分这两种情形所必需的。

为研究一个固定 Einstein 常数 \(\lambda\)，把体积乘子保留在泛函中：
\begin{equation}\label{g3:eq:einstein-action-conserved-tensor}
	\begin{aligned}
		\mathcal A_\lambda(g)&=\mathcal H(g)-(n-2)\lambda V(g),\\
		\mathcal G_\lambda(g)&=\operatorname{Ric}_g-\frac12s_gg+\frac{n-2}{2}\lambda g.
	\end{aligned}
\end{equation}
\(\mathcal G_\lambda\) 是带常数项的 Einstein 张量。它的引入保留了原方程的零点，并使散度恒等式可以直接使用。

\begin{proposition}\label{g3:prop:conserved-einstein-tensor}
	设 \(M^n\) 为闭连通光滑流形，\(n\ge3\)，\(\lambda\in\mathbb R\) 固定。对任意光滑黎曼度量 \(g\)，有
	\begin{equation}\label{g3:eq:einstein-tensor-variation-divergence}
		\mathrm D\mathcal A_\lambda|_g(h)=-\int_M\langle\mathcal G_\lambda(g),h\rangle_g\,\mu_g,
		\qquad \delta_g\mathcal G_\lambda(g)=0.
	\end{equation}
	此外，\(\mathcal G_\lambda(g)=0\) 当且仅当 \(\operatorname{Ric}_g=\lambda g\)。若 \(g_0\) 为这样的 Einstein 度量，记 \(P_0=\mathrm D\mathcal G_\lambda|_{g_0}\)，则
	\begin{equation}\label{g3:eq:einstein-tensor-linearization}
		\begin{aligned}
			P_0h={}&\frac12\Delta_Lh-\delta_{g_0}^*\beta_{g_0}h-\lambda h\\
			&-\frac12\bigl(\Delta_{g_0}\operatorname{tr}_{g_0}h+\delta_{g_0}\delta_{g_0}h-\lambda\operatorname{tr}_{g_0}h\bigr)g_0.
		\end{aligned}
	\end{equation}
	它在 \(L^2(g_0)\) 中形式自伴。
\end{proposition}
\begin{proof}
	泛函微分由上一变分公式直接得到。第3.2节的缩并 Bianchi 恒等式\cref{g3:prop:contracted-bianchi}在当前符号下写成
	\[
		\delta_g\operatorname{Ric}_g=-\frac12\mathrm ds_g.
	\]
	另一方面，\(\delta_g(s_gg)=-\mathrm ds_g\)，且常数 \(\lambda\) 与 \(\nabla g=0\) 给出 \(\delta_g(\lambda g)=0\)。三项相加即得守恒式。

	若 \(\mathcal G_\lambda(g)=0\)，取迹得到
	\[
		0=\operatorname{tr}_g\mathcal G_\lambda(g)=\frac{n-2}{2}(n\lambda-s_g).
	\]
	因为 \(n\ge3\)，故 \(s_g=n\lambda\)；代回即得 \(\operatorname{Ric}_g=\lambda g\)。逆向代入也成立。

	在 Einstein 度量处，微分 \(-s_gg/2+(n-2)\lambda g/2\) 得到 \(-\tfrac12(\mathrm D\operatorname{Scal}_{g_0}h)g_0-\lambda h\)，再代入曲率线性化便得式\eqref{g3:eq:einstein-tensor-linearization}。对两光滑张量 \(h,k\)，在二参数族 \(g_0+th+uk\) 上求 \(\mathcal A_\lambda\) 的混合二阶导数。由于 \(\mathcal G_\lambda(g_0)=0\)，微分积分密度和内积产生的项都为零，所以
	\[
		\mathrm D^2\mathcal A_\lambda|_{g_0}(h,k)=-\int_M\langle P_0h,k\rangle_{g_0}\,\mu_{g_0}.
	\]
	混合导数对称，故右端也对称，这正是 \(P_0\) 的形式自伴性。
\end{proof}

至此有了两种等价的非线性写法：\(\mathcal E_\lambda=0\) 直接表达 Ricci 与度量成比例，\(\mathcal G_\lambda=0\) 则把同一方程放进变分与散度守恒的形式。固定 \(\lambda\) 之后，\(\lambda\ne0\) 时常数缩放会改变方程的参数；\(\lambda=0\) 时缩放仍保留解。下一节采用 \(\mathcal G_\lambda\) 构造规范算子，剩余缩放在有限维核中保留。
% 纲要条目排为编号小节标题，并进入目录与书签；说明段仍保留于纲要注释中。

% !TeX root = ../../Book3.tex
% 纲要标题：### 25.2 规范固定与椭圆化
% 纲要位置：几何学计划纲要.md，第 3176 行。
% 正文状态：初稿；尚未正式审核。

% 纲要内容开始
% 
% * 散度、Bianchi 算子和规范条件
% * 规范化 Einstein 算子的主符号
% * 切片所需的椭圆方程
% * 无穷小等距与剩余对称
% 
% 纲要内容结束

\section{规范固定与椭圆化}
\label{g3:sec:2502}

同一度量经过坐标重参数化，会给出无穷多个张量表达。若直接对 Ricci 方程求逆，这些表达沿微分同胚轨道产生的方向便落在线性化的核中。规范固定要在每条附近的轨道上选一个代表，再把这个代表条件加入方程。这里使用散度条件，并保留 Einstein 张量的散度恒等式；这样既能得到椭圆算子，也能证明新增的规范项在解上消失。

本节固定闭连通光滑流形 \(M^n\)、\(n\ge3\)，以及光滑 Einstein 度量 \(g_0\)，满足 \(\operatorname{Ric}_{g_0}=\lambda g_0\)。下标 \(0\) 表示用 \(g_0\) 定义的微分、内积与体积密度。所有 Sobolev 范数均以这个背景计算。整数 \(s>n/2+3\) 将保证所用张量至少有三阶连续导数。

\subsection{散度、Bianchi 算子和规范条件}
\label{g3:subsec:2502:01}
回顾上一节的约定
\[
	(\delta_g h)_j=-\nabla^i h_{ij},\qquad
	(\delta_g^*\omega)_{ij}=\tfrac12(\nabla_i\omega_j+\nabla_j\omega_i),\qquad
	\beta_g h=\delta_g h+\tfrac12\mathrm d\operatorname{tr}_g h.
\]
闭性使 \(\delta_g^*\) 确为 \(\delta_g\) 的 \(L^2(g)\) 伴随。令 \(X^\flat=g(X,\cdot)\)，则 \(\mathcal L_Xg=2\delta_g^*X^\flat\)。因此散度为零的张量与所有无穷小重参数化正交。我们采用的代表条件为
\begin{equation}\label{g3:eq:metric-divergence-gauge}
	\delta_0(g-g_0)=0.
\end{equation}
这里固定的是背景散度；不能把它换成 \(\delta_g g=0\)，后者对每个度量都自动成立。

\begin{proposition}\label{g3:prop:bianchi-gauge-symbol}
	设 \((M^n,g)\) 为光滑黎曼流形，\(n\ge2\)，\(\omega\) 为光滑一形式。则
	\begin{equation}\label{g3:eq:bianchi-killing-operator}
		2\beta_g\delta_g^*\omega
		=\nabla^*\nabla\omega-\operatorname{Ric}_g(\omega,\cdot).
	\end{equation}
	算子 \(2\delta_g\delta_g^*\) 的主符号为
	\[
		\omega\longmapsto |\xi|_g^2\omega+\xi\langle\xi,\omega\rangle_g,
		\qquad \xi\ne0,
	\]
	因而强椭圆。若 \(M\) 闭，则其核恰为满足 \(\delta_g^*\omega=0\) 的一形式。
\end{proposition}
\begin{proof}
	在一点取正交法标架。由协变导数交换公式，
	\[
		(2\delta_g\delta_g^*\omega)_j
		=(\nabla^*\nabla\omega)_j+\nabla_j\delta_g\omega
		-\operatorname{Ric}_j{}^a\omega_a.
	\]
	又 \(\operatorname{tr}_g\delta_g^*\omega=-\delta_g\omega\)，故加上 \(\mathrm d\operatorname{tr}_g\delta_g^*\omega\) 后中间项抵消，得到式\eqref{g3:eq:bianchi-killing-operator}。用 \(\partial_j\mapsto i\xi_j\) 取最高阶项，得到所列主符号；它与 \(\omega\) 的配对为 \(|\xi|^2|\omega|^2+|\langle\xi,\omega\rangle|^2>0\)。最后在闭流形上
	\[
		\langle\delta_g\delta_g^*\omega,\omega\rangle_{L^2(g)}
		=\|\delta_g^*\omega\|_{L^2(g)}^2,
	\]
	故核的描述成立。
\end{proof}

满足 \(\mathcal L_Xg=0\) 的向量场称为 \(g\) 的\term[Killing chang]{Killing 场}[Killing field]，它生成局部等距变换。记
\[
	\mathfrak k_0=\{\omega\in C^\infty(\mathrm T^*M):\delta_0^*\omega=0\}.
\]
由椭圆核的有限维性，\(\mathfrak k_0\) 有限维。Bianchi 条件 \(\beta_0(g-g_0)=0\) 也常用于 Einstein 方程；式\eqref{g3:eq:bianchi-killing-operator}说明它沿微分同胚方向产生 \(\nabla^*\nabla-\operatorname{Ric}\)。在正 Ricci 曲率下，不能仅凭这个表达式断言其核只有 Killing 场。散度条件给出非负算子 \(\delta_0\delta_0^*\)，核的几何意义则由上面的积分恒等式直接确定。

\begin{example}\label{g3:exm:bianchi-nonkilling-kernel}
	在两份截面曲率为 \(\kappa>0\) 的圆球乘积 \(S^2\times S^2\) 上，乘积度量满足 \(\operatorname{Ric}=\kappa g\)。设 \(f\) 为第一因子上的非零高度函数，拉回到乘积后有 \(\Delta f=2\kappa f\)、\(\operatorname{Hess}f=-\kappa f g_1\)。一形式的 Weitzenböck 公式给出
	\[
		\nabla^*\nabla\mathrm df
		=\mathrm d\Delta f-\operatorname{Ric}(\mathrm df,\cdot)
		=\kappa\,\mathrm df.
	\]
	因此 \(\beta\delta^*\mathrm df=0\)，但 \(\delta^*\mathrm df=\operatorname{Hess}f\ne0\)。这是 Bianchi 规范算子具有非 Killing 核的具体例子。
\end{example}

\subsection{规范化 Einstein 算子的主符号}
\label{g3:subsec:2502:02}
固定 \(\lambda\)，定义
\begin{equation}\label{g3:eq:einstein-conserved-gauge}
	\begin{aligned}
		\mathcal G_\lambda(g)&=\operatorname{Ric}_g-\tfrac12s_g g
			+\tfrac{n-2}{2}\lambda g,\\
		\omega(g)&=\delta_0(g-g_0),\\
		\mathcal F(g)&=\mathcal G_\lambda(g)+\delta_g^*\omega(g).
	\end{aligned}
\end{equation}
缩并 Bianchi 恒等式给出 \(\delta_g\mathcal G_\lambda(g)=0\)。另一方面，\(\mathcal G_\lambda(g)=0\) 的迹为 \((n-2)(n\lambda-s_g)/2=0\)，所以在 \(n\ge3\) 时它等价于 \(\operatorname{Ric}_g=\lambda g\)。以这个守恒张量代替 Ricci 张量，正是随后消去规范误差的原因。

\begin{theorem}\label{g3:thm:einstein-gauged-ellipticity}
	设 \(M^n\) 闭且光滑，\(n\ge3\)，\(g_0\) 为满足 \(\operatorname{Ric}_{g_0}=\lambda g_0\) 的光滑度量。取整数 \(s>n/2+3\)。式\eqref{g3:eq:einstein-conserved-gauge}定义光滑映射
	\[
		\mathcal F:\mathcal M^s(M)\longrightarrow
		H^{s-2}(S^2\mathrm T^*M).
	\]
	在 \(g_0\) 附近它是二阶准线性椭圆系统。其线性化
	\begin{equation}\label{g3:eq:einstein-gauged-linearization}
		Lh=\mathrm D\mathcal G_\lambda|_{g_0}h+\delta_0^*\delta_0h
	\end{equation}
	关于 \(L^2(g_0)\) 形式自伴。
\end{theorem}
\begin{proof}
	局部表达式中，逆度量是正定矩阵的光滑函数，曲率为 \(g^{-1}*\partial^2g+g^{-1}*g^{-1}*\partial g*\partial g\) 的有限和。规范项也至多包含二阶导数。\Cref{g3:prop:nonlinear-jet-operator}适用，因为 \(s-2>n/2\)，故映射光滑。由于 \(\omega(g_0)=0\)，微分规范项时只有 \(\delta_0^*\delta_0h\) 留下。

	令 \(\xi\) 为单位余向量，在 \(g_0\) 正交标架中取 \(\xi=e^1\)。上一节的 Ricci 线性化给出
	\[
		\sigma_{\operatorname{Ric}'}(\xi)h
		=\tfrac12\bigl(h+\xi\otimes\xi\operatorname{tr}h
		-\xi\otimes h(\xi^\sharp,\cdot)-h(\xi^\sharp,\cdot)\otimes\xi\bigr),
	\]
	而标量曲率的符号是 \(\operatorname{tr}h-h(\xi^\sharp,\xi^\sharp)\)。加上 \(\delta_0^*\delta_0\) 后，交叉项抵消。若 \(a,b\ge2\)、\(H=\sum_{c=2}^nh_{cc}\)，则
	\begin{equation}\label{g3:eq:einstein-gauge-block-symbol}
		(\sigma_Lh)_{11}=h_{11},\quad
		(\sigma_Lh)_{1a}=\tfrac12h_{1a},\quad
		(\sigma_Lh)_{ab}=\tfrac12(h_{ab}-H\delta_{ab}).
	\end{equation}
	最后一块的迹是 \(-(n-2)H/2\)。因此 \(\sigma_Lh=0\) 依次推出 \(h_{11}=h_{1a}=H=h_{ab}=0\)。任意非零 \(\xi\) 由齐次性得到可逆性。单位余球丛紧，最小奇异值有正下界；系数连续依赖 \(g\)，故附近度量仍椭圆。这个符号有负的迹方向，椭圆性并不等于能量正定。

	上一节的作用量 \(\mathcal A_\lambda(g)=\int_Ms_g\,\mu_g-(n-2)\lambda\operatorname{Vol}_g(M)\) 满足
	\(\mathrm D\mathcal A_\lambda(g)h=-\langle\mathcal G_\lambda(g),h\rangle_{L^2(g)}\)。在 \(g_0\) 处 \(\mathcal G_\lambda=0\)，故再微分时内积和体积的变化项都消失，Hessian 的对称性给出 \(\mathcal G_\lambda'\) 的形式自伴性。算子 \(\delta_0^*\delta_0\) 同样自伴，得到结论。
\end{proof}

\begin{theorem}\label{g3:thm:einstein-gauge-equivalence}
	设 \((M^n,g_0)\) 为闭光滑 Einstein 流形，\(n\ge3\)，\(\operatorname{Ric}_{g_0}=\lambda g_0\)。记 \(\delta_0=\delta_{g_0}\) 为负散度、\(\delta_g^*\) 为其对称梯度伴随，并令 \(\mathcal F(g)=\operatorname{Ric}_g-\tfrac12s_gg+\tfrac{n-2}2\lambda g+\delta_g^*\delta_0(g-g_0)\)。存在 \(g_0\) 的一个 \(C^1\) 邻域，使其中每个 \(C^3\) 度量 \(g\) 满足
	\[
		\mathcal F(g)=0
		\quad\Longleftrightarrow\quad
		\operatorname{Ric}_g=\lambda g\ \text{且}\ \delta_0(g-g_0)=0.
	\]
\end{theorem}
\begin{proof}
	只需证明正向。设 \(\omega=\delta_0(g-g_0)\)。对方程取 \(\delta_g\)，由守恒恒等式得到 \(\delta_g\delta_g^*\omega=0\)。与 \(\omega\) 配对并积分，得 \(\delta_g^*\omega=0\)，于是 \(\mathcal G_\lambda(g)=0\)，亦即 \(\operatorname{Ric}_g=\lambda g\)。

	还须证明代表条件本身。对任意 \(\eta\in\mathfrak k_0\)，有
	\(\langle\omega,\eta\rangle_0=\langle g-g_0,\delta_0^*\eta\rangle_0=0\)。在 \(\mathfrak k_0^\perp\) 上有估计
	\begin{equation}\label{g3:eq:killing-complement-first-order}
		\|\alpha\|_{H^1}\le C\|\delta_0^*\alpha\|_2.
	\end{equation}
	为证明它，先对具有单射符号的 \(\delta_0^*\) 冻结系数、用 Fourier 估计及有限覆盖，得 \(\|\alpha\|_{H^1}\le C(\|\delta_0^*\alpha\|_2+\|\alpha\|_2)\)。若不能去掉末项，取 \(\mathfrak k_0^\perp\) 中 \(H^1\) 范数为一而 \(\delta_0^*\alpha_j\to0\) 的序列。Rellich 紧性给出强 \(L^2\) 极限 \(\alpha\)；弱方程使它属于 \(\mathfrak k_0\)，正交性又使它为零。前述估计遂给 \(\alpha_j\to0\) 于 \(H^1\)，矛盾。

	两联络之差是由 \(g^{-1}\nabla^0g\) 给出的零阶张量，因此
	\[
		\|\delta_0^*\omega\|_2
		=\|(\delta_0^*-\delta_g^*)\omega\|_2
		\le C'\|g-g_0\|_{C^1}\|\omega\|_2.
	\]
	取邻域使 \(CC'\|g-g_0\|_{C^1}<1\)，式\eqref{g3:eq:killing-complement-first-order}迫使 \(\omega=0\)。反向代入定义即得。
\end{proof}

\subsection{切片所需的椭圆方程}
\label{g3:subsec:2502:03}
先把线性代表条件解出来。给定对称二张量 \(h\)，希望沿 \(h\mapsto h+\mathcal L_Xg_0\) 消去散度，因此要解
\begin{equation}\label{g3:eq:metric-linear-slice-equation}
	2\delta_0\delta_0^*X^\flat=-\delta_0h.
\end{equation}
右端自动与 \(\mathfrak k_0\) 正交。\Cref{g3:thm:elliptic-fredholm,g3:cor:elliptic-range-hk}因而给出唯一的核正交解，且 \(\|X\|_{H^{r+1}}\le C_r\|h\|_{H^r}\)。这也给出正交分解
\[
	H^r(S^2\mathrm T^*M)
	=\ker\delta_0\ \oplus\ \delta_0^*H^{r+1}(\mathrm T^*M),\qquad r\ge1.
\]
第二项不要求把 Killing 场重复计入参数。

非线性求解涉及变化源上的复合。为看清正则性，令
\(\phi_X(x)=\exp^{g_0}_xX(x)\)，其中 \(X\) 足够 \(C^1\) 小。它是邻近恒等的微分同胚：局部逆函数定理给局部可逆，紧性给覆盖，而与恒等的同伦使覆盖次数为一。

\begin{lemma}\label{g3:lem:metric-pullback-one-loss}
	设 \(M^n\) 闭且光滑，\(g_0\) 为光滑黎曼度量，整数 \(s>n/2+3\)。记 \(\phi_X(x)=\exp_x^{g_0}X(x)\)，\(\delta_0^*\) 为背景度量的对称梯度，\(\mathcal M^{s+1}(M)\) 为正定 \(H^{s+1}\) 对称二张量的开集。在 \(X=0\)、\(g=g_0\) 附近，映射
	\[
		(X,g)\longmapsto\phi_X^*g
	\]
	从 \(H^{s+1}(\mathrm TM)\times\mathcal M^{s+1}(M)\) 到 \(H^s(S^2\mathrm T^*M)\) 连续可微；其在 \((0,g_0)\) 处的微分是 \((X,h)\mapsto2\delta_0^*X^\flat+h\)。
\end{lemma}
\begin{proof}
	有限坐标覆盖上只需检查
	\((\phi^*g)_{ij}(x)=g_{ab}(\phi(x))\partial_i\phi^a\partial_j\phi^b\)。先说明复合估计：在 \(H^{s+1}\) 有界且 Jacobian 及其逆有界的微分同胚邻域中，\(f\mapsto f\circ\phi\) 在每个整数 \(0\le r\le s+1\) 的 \(H^r\) 上有统一界。\(r=0\) 由换元得出；从 \(r-1\) 到 \(r\) 对复合求一次导数，用 \(\mathrm D(f\circ\phi)=(\mathrm Df\circ\phi)\mathrm D\phi\) 及\cref{g3:lem:sobolev-multiple-products}的乘积估计归纳。对光滑 \(f\) 有关于 \(\phi\) 的连续性，一般 \(f\) 由光滑稠密性和统一界逼近，因此联合连续。

	局部坐标内，\(\phi\) 的增量 \(V\) 给出
	\[
		(f\circ(\phi+V)-f\circ\phi)
		=\int_0^1(\mathrm Df\circ(\phi+tV))V\,\mathrm dt.
	\]
	减去 \((\mathrm Df\circ\phi)V\)，联合连续性说明余项是 \(o(\|V\|_{H^{s+1}})\) 于 \(H^s\)。同一积分式与统一复合界还给出局部估计
	\[
		\|f\circ\phi-f\circ\psi\|_{H^s}
		\le C\|f\|_{H^{s+1}}\|\phi-\psi\|_{H^{s+1}}.
	\]
	它保证关于 \(f\) 的微分在算子范数中连续，并控制 \(f\) 与 \(\phi\) 同时变化时的交叉项；关于 \(\phi\) 的微分则由 \(\mathrm Df\circ\phi\) 在 \(H^s\) 中的连续性控制。故复合从两份 \(H^{s+1}\) 到 \(H^s\) 为 \(C^1\)。拉回表达式中的两个 \(\mathrm D\phi\) 属于 \(H^s\)，而 \(H^s\) 是乘法代数；代入纤维指数图，链式法则给出结论。微分的坐标表达式为 Lie 导数，故得到所列公式。
\end{proof}

以下正则性事实也将用于下一节的有限维约化。需要的是一般等秩椭圆系统，而不是只含标量主部的方程。

\begin{lemma}\label{g3:lem:elliptic-system-holder-bootstrap}
	设 \(u\) 是开集上的 \(C^{3,\alpha}\) 向量值函数，\(0<\alpha<1\)，满足
	\[
		A^{ij}(x,u,\mathrm Du)\partial_i\partial_ju
		=B(x,u,\mathrm Du)+f(x),
	\]
	其中 \(A,B,f\) 光滑，主符号 \(\sum A^{ij}(x,u,\mathrm Du)\xi_i\xi_j\) 在每个紧子集和单位 \(\xi\) 上可逆且逆一致有界。则 \(u\) 光滑。若主系数只依赖 \(x,u\)，结论对相同起始正则性同样成立；在固定紧子集上，每阶估计由相应低阶范数、椭圆性和数据界控制。
\end{lemma}
\begin{proof}
	先把第12.5节的内部 Schauder 证明扩展到矩阵主符号。设冻结算子为 \(P_0\)，其二阶符号为 \(p_2(\xi)\)。对二阶导数求解所用的乘子是 \(-\xi_k\xi_l p_2(\xi)^{-1}\)，它是零次齐次、在单位球面光滑的矩阵函数。分离常矩阵乘法项后，\cref{g3:thm:smooth-multiplier-lp}证明中的二进分解给出齐次核 \(K\) 的界 \(|K(x)|\le C|x|^{-n}\)、\(|\mathrm DK(x)|\le C|x|^{-n-1}\)，以及球面零均值所给的环域消去。对 \(f\in C^{0,\alpha}\)，在近区减去 \(f(x)\)，积分界为 \(C[f]_\alpha\int_0^r t^{\alpha-1}\,\mathrm dt\)；比较相距 \(r\) 的两点时，远区用核差，界为 \(Cr[f]_\alpha\int_r^1t^{\alpha-2}\,\mathrm dt\)。两项均为 \(O(r^\alpha)\)。取截断函数 \(\chi\)，把上述全空间估计用于 \(\chi v\)，利用 \(P_0(\chi v)=\chi P_0v+[P_0,\chi]v\)；交换子只含一阶导数，可用加权插值吸收。这就给出常系数系统的内部二阶 Hölder 估计。

	对变系数系统，冻结误差仍是矩阵乘积 \((A(x)-A(x_0))\mathrm D^2v\)。缩小球吸收最高阶半范数，并用\cref{g3:lem:weighted-holder-interpolation}吸收其余导数项，得到内部估计
	\[
		\|v\|_{C^{2,\alpha}(U')}
		\le C\bigl(\|Av\|_{C^{0,\alpha}(U)}+\|v\|_{C^0(U)}\bigr),
		\qquad U'\Subset U.
	\]
	此处只使用矩阵范数和逆符号界。对初始 \(C^2\) 解，\cref{g3:eq:holder-mollifier-commutator}的卷积交换子估计逐矩阵项适用，故估计包含正则性结论。对差商使用这个估计，正如\cref{g3:cor:schauder-higher-bootstrap}的证明，系数和右端为 \(C^{r,\alpha}\) 时得到 \(C^{r+2,\alpha}\) 正则性。

	回到本引理，初始 \(u\in C^{3,\alpha}\) 使 \(A(x,u,\mathrm Du)\) 和右端至少属于 \(C^{2,\alpha}\)。将这些已经确定的函数作为线性方程的系数，得到 \(u\in C^{4,\alpha}\)。系数和右端随之成为 \(C^{3,\alpha}\)，再得 \(C^{5,\alpha}\)。归纳完成证明。每一步的估计也给出所述依赖关系。
\end{proof}

\begin{theorem}\label{g3:thm:metric-divergence-slice}
	设 \(M^n\) 闭且光滑，\(g_0\) 为光滑黎曼度量，整数 \(s>n/2+3\)。记 \(\delta_0=\delta_{g_0}\)、\(\mathfrak k_0=\ker\delta_0^*\)，正交取 \(L^2(g_0)\) 内积；对充分小的 \(X\)，令 \(\phi_X(x)=\exp_x^{g_0}X(x)\)。每个充分接近 \(g_0\) 的 \(H^{s+1}\) 度量 \(g\)，均有唯一充分小的 \(X\in H^{s+1}(\mathrm TM)\)、\(X^\flat\perp\mathfrak k_0\)，使
	\[
		\delta_0(\phi_X^*g-g_0)=0.
	\]
	此时 \(\phi_X^*g\in H^s\)，\(X\) 对 \(g\) 连续可微。若 \(g\) 光滑，则 \(X\)、\(\phi_X\) 及所得代表都光滑。
\end{theorem}
\begin{proof}
	令 \(\Phi(X,g)=\delta_0(\phi_X^*g-g_0)\)。前一复合引理说明它取值于 \(H^{s-1}\) 且为 \(C^1\)。其像始终正交于 \(\mathfrak k_0\)，因为背景散度的像具有这个性质。在 \((0,g_0)\) 处，\(X\) 微分是 \(2\delta_0\delta_0^*X^\flat\)。\Cref{g3:prop:bianchi-gauge-symbol}与椭圆 Fredholm 选择表明它是
	\[
		H^{s+1}(\mathrm T^*M)\cap\mathfrak k_0^\perp
		\longrightarrow H^{s-1}(\mathrm T^*M)\cap\mathfrak k_0^\perp
	\]
	之间的有界同构。应用\cref{g3:cor:banach-implicit-parameters}，得到存在、唯一性及 \(C^1\) 依赖。

	若 \(g\) 光滑，写出 \(\Phi(X,g)=0\) 的坐标表达式，它是关于 \(X\) 的二阶准线性系统；最高阶系数光滑依赖 \(x,X,\mathrm DX\)，在 \(X=0,g=g_0\) 处为上述强椭圆符号。邻域缩小后保持椭圆。由 \(X\in H^{s+1}\) 得 \(X\in C^{3,\alpha}\) 对某个 \(\alpha>0\) 成立，故\cref{g3:lem:elliptic-system-holder-bootstrap}给出光滑性。
\end{proof}

这个证明把拉回的一阶微分损失明确放进了空间选择：输入度量为 \(H^{s+1}\)，所得代表先在 \(H^s\) 中控制。它没有把微分同胚作用当作同一 Sobolev 阶上的光滑作用；光滑度量的几何结论则由椭圆正则性恢复。

\subsection{无穷小等距与剩余对称}
\label{g3:subsec:2502:04}
令 \(I_0=\operatorname{Isom}(M,g_0)\)。它保持散度条件，也保持规范化方程：对 \(a\in I_0\)，
\[
	\delta_0(a^*h)=a^*(\delta_0h),\qquad
	\mathcal F(a^*g)=a^*\mathcal F(g).
\]
因而规范固定后仍须除去 \(I_0\) 的作用。其 Lie 代数为 Killing 场，但离散等距变换同样可能起作用。

\begin{proposition}\label{g3:prop:metric-slice-residual-isometries}
	设 \(M\) 闭且连通，\(g_0\) 光滑。记 \(\delta_0=\delta_{g_0}\) 为负散度，\(I_0=\operatorname{Isom}(M,g_0)\)。光滑度量在 \(g_0\) 附近满足 \(\delta_0(g-g_0)=0\) 时，模所有微分同胚的局部等价关系恰由 \(I_0\) 的拉回作用给出。这里“附近”取光滑拓扑，并可相应缩小上一切片定理的邻域。
\end{proposition}
\begin{proof}
	先考察恒等附近的微分同胚。Killing 场的流保持 \(g_0\)，由闭性对所有小时间存在。选 \(\mathfrak k_0\) 的一组基，按次序复合这些流，得到有限维参数 \(a(t)\)。映射 \((X,t)\mapsto a(t)\circ\phi_X\)，其中 \(X^\flat\perp\mathfrak k_0\)，在 \((0,0)\) 的微分是核正交向量场与 Killing 场之和，故为同构。这里外层流光滑且仅有有限个参数，映射在 \(H^{s+1}\) 微分同胚图中光滑。逆函数定理说明每个充分靠近恒等的微分同胚都有这种表示。

	若 \(g_1,g_2\) 均在散度切片内，且 \(g_2=(a(t)\circ\phi_X)^*g_1=\phi_X^*(a(t)^*g_1)\)，则 \(a(t)^*g_1\) 也在切片内。切片定理对这一输入的唯一解是 \(X=0\)，因此该等价来自等距变换。

	最后排除远离 \(I_0\) 的额外等价。若有 \(g_j,h_j\to g_0\) 光滑，\(\psi_j^*g_j=h_j\)，则拉回恒等式给出 \(\mathrm D\psi_j\) 及其逆的一致界。局部等距映射保持联络，因此在坐标中
	\[
		\partial_i\partial_j\psi_j^a
		=\Gamma(h_j)^k_{ij}\partial_k\psi_j^a
		-\Gamma(g_j)^a_{bc}(\psi_j)\partial_i\psi_j^b\partial_j\psi_j^c.
	\]
	有限坐标覆盖与此式依次控制全部高阶导数；逆映射有相同估计。由 Arzelà--Ascoli 及对角取列，\(\psi_j\) 和逆映射光滑收敛到互逆的映射 \(\psi,\psi^{-1}\)，并且 \(\psi^*g_0=g_0\)。故 \(\psi\in I_0\)，\(\psi^{-1}\circ\psi_j\to\operatorname{id}\)。应用上一段，便得到所需结论；否则一列违反结论的等价会产生矛盾。
\end{proof}

同一紧性论证表明 \(I_0\) 紧。恒等附近的分解还说明其中靠近恒等的元素全部由 Killing 场的有限维流参数给出：对 \(g_1=g_2=g_0\) 再用切片唯一性即可。于是 \(I_0\) 是有限维紧 Lie 群；这也说明为何它的离散部分不能由核空间计算替代。上面的论证是本章使用的散度切片版本；度量空间切片的一般框架见\cite[第7节，定理7.4，第33--34页]{Ebin1970Metrics}。
% 纲要条目排为编号小节标题，并进入目录与书签；说明段仍保留于纲要注释中。

% !TeX root = ../../Book3.tex
% 纲要标题：### 25.3 无穷小形变及障碍
% 纲要位置：几何学计划纲要.md，第 3183 行。
% 正文状态：初稿；尚未正式审核。

% 纲要内容开始
% 
% * Lichnerowicz 型算子
% * 核、余核和规范等价
% * 隐函数定理的可解条件
% * 有限维约化及二阶障碍
% 
% 纲要内容结束

\section{无穷小形变及障碍}
\label{g3:sec:2503}

椭圆化把无穷多个重参数化方向移出了核，却未保证剩余的线性化可逆。真正的几何形变可能留下有限维核，而核中的向量未必能够延伸为一条 Einstein 度量曲线。本节先识别这些无穷小方向，再解出核正交部分，把全部非线性困难保存在一个有限维方程中。

固定闭连通光滑 Einstein 流形 \((M^n,g_0)\)，\(n\ge3\)，\(\operatorname{Ric}_{g_0}=\lambda g_0\)，以及整数 \(s>n/2+3\)。本节先固定 Einstein 常数 \(\lambda\)；最后再将它作为一个实参数加入，以描述未固定尺度的局部问题。

\subsection{Lichnerowicz 型算子}
\label{g3:subsec:2503:01}
对光滑对称二张量 \(h\)，记 \(f=\operatorname{tr}_0h\)。上一节的线性化可以写为
\begin{equation}\label{g3:eq:einstein-L-expanded}
	Lh=\tfrac12(\Delta_L-2\lambda)h
	-\tfrac12\operatorname{Hess}_0 f
	-\tfrac12\bigl(\Delta_0f+\delta_0\delta_0h-\lambda f\bigr)g_0.
\end{equation}
这里 \(\Delta_L\) 是第一节定义的 Lichnerowicz 算子；其在 Einstein 背景上满足
\[
	(\Delta_L-2\lambda)h=\nabla^{0*}\nabla^0h-2\mathring R_0(h).
\]
式\eqref{g3:eq:einstein-L-expanded}来自 \(\operatorname{Ric}'=\Delta_L/2-\delta_0^*\beta_0\)：规范项消去 \(-\delta_0^*\delta_0h\)，而 Einstein 张量的迹反转留下最后一项。不能只保留粗 Laplace 项而忽略这些迹项。

\begin{definition}\label{g3:def:einstein-infinitesimal-space}
	设 \((M^n,g_0)\) 为闭光滑 Einstein 流形，\(n\ge3\)，\(\operatorname{Ric}_{g_0}=\lambda g_0\)。对背景 Einstein 度量 \(g_0\)，以 \(\delta_0\) 表示负散度，令 \(\mathcal F(g)=\operatorname{Ric}_g-\tfrac12s_gg+\tfrac{n-2}2\lambda g+\delta_g^*\delta_0(g-g_0)\)，\(L=\mathrm D\mathcal F|_{g_0}\)，\(\mathcal K=\ker L\) 取光滑对称二张量核；所有正交与内积取 \(L^2(g_0)\)。记
	\[
		\mathcal K=\ker L\subset C^\infty(S^2\mathrm T^*M).
	\]
	满足 \(\delta_0h=0\)、\(\operatorname{tr}_0h=0\) 的张量称为\term[hengxiang wuji zhangliang]{横向无迹张量}[transverse-traceless tensor]，简称 TT 张量。空间 \(\mathcal K\) 中的元素称为固定 \(\lambda\) 的规范\term[wuqiongxiao Einstein xingbian]{无穷小 Einstein 形变}[infinitesimal Einstein deformation]。
\end{definition}
\symidx{\(\mathcal K\)：规范固定 Einstein 方程线性化算子的光滑核}

\begin{proposition}\label{g3:prop:einstein-kernel-tt}
	设 \((M^n,g_0)\) 闭、连通且光滑，\(n\ge3\)，\(\operatorname{Ric}_{g_0}=\lambda g_0\)。对背景 Einstein 度量 \(g_0\)，以 \(\delta_0\) 表示负散度，令 \(\mathcal F(g)=\operatorname{Ric}_g-\tfrac12s_gg+\tfrac{n-2}2\lambda g+\delta_g^*\delta_0(g-g_0)\)，\(L=\mathrm D\mathcal F|_{g_0}\)，\(\mathcal K=\ker L\) 取光滑对称二张量核；所有正交与内积取 \(L^2(g_0)\)。则
	\begin{equation}\label{g3:eq:einstein-kernel-geometric}
		\mathcal K=\{h:\delta_0h=0,\ \mathrm D(\operatorname{Ric}-\lambda g)|_{g_0}h=0\}.
	\end{equation}
	若 \(\lambda\ne0\)，则
	\[
		\mathcal K=\{h:\delta_0h=0,\ \operatorname{tr}_0h=0,
		(\nabla^{0*}\nabla^0-2\mathring R_0)h=0\}.
	\]
	若 \(\lambda=0\)，则 \(\mathcal K=\mathbb Rg_0\oplus\mathcal K_{\mathrm{TT}}\)，其中 \(\mathcal K_{\mathrm{TT}}\) 为上式中的 TT 核。
\end{proposition}
\begin{proof}
	微分 \(\delta_g\mathcal G_\lambda(g)=0\)，并用 \(\mathcal G_\lambda(g_0)=0\)，得到 \(\delta_0\mathcal G_\lambda'h=0\)。若 \(Lh=0\)，则 \(\delta_0\delta_0^*\delta_0h=0\)。将它与 \(\delta_0h\) 配对，得 \(\delta_0^*\delta_0h=0\)；又 \(\delta_0h\perp\mathfrak k_0\)，故 \(\delta_0h=0\)。于是 \(\mathcal G_\lambda'h=0\)。令 \(E'h=\mathrm D(\operatorname{Ric}-\lambda g)|_{g_0}h\)，由第一节的变分公式有
	\[
		\mathcal G_\lambda'h=E'h-\tfrac12(\operatorname{tr}_0E'h)g_0.
	\]
	迹反转在 \(n\ge3\) 时可逆，故 \(E'h=0\)。反向显然代入即得，证明式\eqref{g3:eq:einstein-kernel-geometric}。

	对这个方程取迹，标量曲率变分给出 \(\Delta_0f=\lambda f\)。若 \(\lambda<0\)，积分恒等式 \(\|\mathrm df\|_2^2=\lambda\|f\|_2^2\) 迫使 \(f=0\)。若 \(\lambda>0\)，先积分得 \(\int_M f=0\)。\Cref{g3:thm:lichnerowicz-bound}在 \(\operatorname{Ric}_{g_0}=\lambda g_0\) 下给出首个正函数特征值至少为 \(n\lambda/(n-1)>\lambda\)，故仍有 \(f=0\)。此时 \(\beta_0h=0\)，于是 \(E'h=(\nabla^{0*}\nabla^0-2\mathring R_0)h/2\)。

	若 \(\lambda=0\)，则 \(\Delta_0f=0\)，闭性与连通性给出 \(f\) 为常数。平行张量 \(g_0\) 满足 \(L g_0=0\)；从 \(h\) 减去 \((f/n)g_0\) 就得到 TT 核。两项正交，给出直和分解。
\end{proof}

因此尺度的地位取决于规范化：非零 \(\lambda\) 已经固定尺度，而 Ricci 平坦方程在常数缩放下不变。TT 条件描述的是去掉这一尺度方向后的变化。

\subsection{核、余核和规范等价}
\label{g3:subsec:2503:02}
核记录齐次线性方程的解，余核记录非齐次方程的相容条件。两者在当前自伴问题中可用同一个光滑有限维空间表示，但在非线性约化中仍承担不同职责。

\begin{theorem}\label{g3:thm:einstein-linear-fredholm}
	设 \((M^n,g_0)\) 为闭光滑 Einstein 流形，\(n\ge3\)，记 \(\operatorname{Ric}_{g_0}=\lambda g_0\)，取整数 \(s>n/2+3\)。对背景 Einstein 度量 \(g_0\)，以 \(\delta_0\) 表示负散度，令 \(\mathcal F(g)=\operatorname{Ric}_g-\tfrac12s_gg+\tfrac{n-2}2\lambda g+\delta_g^*\delta_0(g-g_0)\)，\(L=\mathrm D\mathcal F|_{g_0}\)，\(\mathcal K=\ker L\) 取光滑对称二张量核；所有正交与内积取 \(L^2(g_0)\)。则 \(\mathcal K\) 有限维，且
	\[
		L:H^s(S^2\mathrm T^*M)\longrightarrow H^{s-2}(S^2\mathrm T^*M)
	\]
	为指数零的 Fredholm 算子。对 \(q\in H^{s-2}\)，方程 \(Lh=q\) 可解当且仅当 \(\langle q,k\rangle_0=0\) 对所有 \(k\in\mathcal K\) 成立；核正交解唯一，并满足
	\begin{equation}\label{g3:eq:einstein-complement-estimate}
		\|h\|_{H^s}\le C_s\|Lh\|_{H^{s-2}},\qquad h\perp\mathcal K.
	\end{equation}
\end{theorem}
\begin{proof}
	\Cref{g3:thm:einstein-gauged-ellipticity}给出等秩二阶椭圆性及形式自伴性。\Cref{g3:prop:finite-elliptic-kernel,g3:cor:elliptic-range-hk}遂给出有限维光滑核、闭像以及像等于伴随核的正交补。这里伴随算子仍为 \(L\)，故余核与 \(\mathcal K\) 同维，指数为零。核正交估计由\cref{g3:thm:kernel-orthogonal-estimate}取阶数二、指数二得到。注意这一论证只需椭圆性，不要求 \(L\) 的二次型非负。
\end{proof}

取 \(\mathcal K\) 的 \(L^2(g_0)\) 正交归一基 \(k_1,\ldots,k_d\)，定义
\[
	\Pi h=\sum_{a=1}^d\langle h,k_a\rangle_0k_a,\qquad Q=I-\Pi.
\]
各 \(k_a\) 光滑，因此 \(\Pi\) 在本节所用的两个 Sobolev 空间上都连续。令
\[
	X_1=H^s\cap\mathcal K^\perp,\qquad Y_1=H^{s-2}\cap\mathcal K^\perp.
\]
则 \(L:X_1\to Y_1\) 是有界同构，逆记为 \(G\)。它是固定背景下的 Green 逆，与第13.3节的构造一致。

\begin{proposition}\label{g3:prop:einstein-gauge-quotient-tangent}
	设 \((M^n,g_0)\) 为闭光滑 Einstein 流形，\(n\ge3\)，\(\operatorname{Ric}_{g_0}=\lambda g_0\)。对背景 Einstein 度量 \(g_0\)，以 \(\delta_0\) 表示负散度，令 \(\mathcal F(g)=\operatorname{Ric}_g-\tfrac12s_gg+\tfrac{n-2}2\lambda g+\delta_g^*\delta_0(g-g_0)\)，\(L=\mathrm D\mathcal F|_{g_0}\)，\(\mathcal K=\ker L\) 取光滑对称二张量核；所有正交与内积取 \(L^2(g_0)\)。令
	\[
		\mathcal Z=\ker\mathrm D(\operatorname{Ric}-\lambda g)|_{g_0}.
	\]
	则取散度规范给出自然线性同构
	\[
		\mathcal Z/\{\mathcal L_Xg_0:X\in C^\infty(\mathrm TM)\}
		\simeq\mathcal K.
	\]
\end{proposition}
\begin{proof}
	自然性使 \(\mathcal L_Xg_0\in\mathcal Z\)。对 \(h\in\mathcal Z\)，解式\eqref{g3:eq:metric-linear-slice-equation}，得到 \(\widehat h=h+\mathcal L_Xg_0\) 满足 \(\delta_0\widehat h=0\)。它仍在 \(\mathcal Z\)，故属于 \(\mathcal K\)。若一个 Lie 导数本身散度为零，则 \(\delta_0\delta_0^*X^\flat=0\)，积分后得到 \(\delta_0^*X^\flat=0\)，所以该 Lie 导数就是零。因而所选代表唯一。
\end{proof}

这个线性商尚不是实际模空间的切空间定理：实际解集可能有高阶障碍。它首先给出所有可能曲线速度所在的有限维候选空间。

\subsection{隐函数定理的可解条件}
\label{g3:subsec:2503:03}
若 \(\mathcal K=0\)，线性化可直接求逆；若 \(\mathcal K\ne0\)，可逆的是核正交部分。两种情形分别决定隐函数定理的使用方式。

\begin{corollary}\label{g3:cor:einstein-infinitesimal-rigidity}
	设 \((M^n,g_0)\) 为闭光滑 Einstein 流形，\(n\ge3\)，\(\operatorname{Ric}_{g_0}=\lambda g_0\)。对背景 Einstein 度量 \(g_0\)，以 \(\delta_0\) 表示负散度，令 \(\mathcal F(g)=\operatorname{Ric}_g-\tfrac12s_gg+\tfrac{n-2}2\lambda g+\delta_g^*\delta_0(g-g_0)\)，\(L=\mathrm D\mathcal F|_{g_0}\)，\(\mathcal K=\ker L\) 取光滑对称二张量核；所有正交与内积取 \(L^2(g_0)\)。若 \(\mathcal K=0\)，则每个具有同一 Einstein 常数 \(\lambda\)、且充分光滑接近 \(g_0\) 的 Einstein 度量，都与 \(g_0\) 微分同胚等价。
\end{corollary}
\begin{proof}
	\Cref{g3:thm:einstein-linear-fredholm}使 \(L:H^s\to H^{s-2}\) 为同构。\Cref{g3:thm:quantitative-banach-inverse}说明 \(\mathcal F\) 在 \(g_0\) 附近的唯一零点为 \(g_0\)。对邻近的同常数 Einstein 度量，先用\cref{g3:thm:metric-divergence-slice}选散度代表。代表保持 Einstein 方程，由\cref{g3:thm:einstein-gauge-equivalence}是 \(\mathcal F\) 的零点，故等于 \(g_0\)。
\end{proof}

在 \(\lambda=0\) 时 \(g_0\in\mathcal K\)，所以上述条件不会自动满足；需固定体积后再讨论刚性。在一般 Banach 方程中，满射而非零核的线性化会产生光滑解流形；当前 \(L\) 自伴且指数零，满射已经蕴含核为零。若核不为零，不能把“有限维”当作“没有相容条件”。

把 \(h=k+v\)，其中 \(k\in\mathcal K\)、\(v\in X_1\)。首先只解
\begin{equation}\label{g3:eq:einstein-range-equation}
	Q\mathcal F(g_0+k+v)=0.
\end{equation}
它对 \(v\) 在原点的微分是 \(L:X_1\to Y_1\)，所以确实满足隐函数定理。问题中是否存在给定拓扑上的 Einstein 度量，并未由这个局部求逆解决；这里始终从已知解 \(g_0\) 出发。

\subsection{有限维约化及二阶障碍}
\label{g3:subsec:2503:04}
\begin{theorem}\label{g3:thm:einstein-local-obstruction}
	设 \((M^n,g_0)\) 为闭连通光滑 Einstein 流形，\(n\ge3\)，\(\operatorname{Ric}_{g_0}=\lambda g_0\)，取整数 \(s>n/2+3\)。对背景 Einstein 度量 \(g_0\)，以 \(\delta_0\) 表示负散度，令 \(\mathcal F(g)=\operatorname{Ric}_g-\tfrac12s_gg+\tfrac{n-2}2\lambda g+\delta_g^*\delta_0(g-g_0)\)，\(L=\mathrm D\mathcal F|_{g_0}\)，\(\mathcal K=\ker L\) 取光滑对称二张量核；所有正交与内积取 \(L^2(g_0)\)。记 \(X_1=H^s(S^2\mathrm T^*M)\cap\mathcal K^\perp\)。存在 \(\mathcal K\) 中零点的邻域 \(U\) 及光滑映射
	\[
		\psi:U\to X_1,\qquad \kappa:U\to\mathcal K,
	\]
	使 \(\psi(0)=0\)、\(\mathrm D\psi(0)=0\)、\(\kappa(0)=0\)、\(\mathrm D\kappa(0)=0\)，并有：
	\begin{enumerate}
		\item \(g_0\) 附近 \(H^s\) 中的 \(\mathcal F\) 零点恰为 \(g(k)=g_0+k+\psi(k)\)，其中 \(k\in U\)、\(\kappa(k)=0\)。这些度量都是光滑的，满足 \(\operatorname{Ric}_{g(k)}=\lambda g(k)\) 及 \(\delta_0(g(k)-g_0)=0\)。
		\item 每个充分光滑接近 \(g_0\)、具有同一 \(\lambda\) 的光滑 Einstein 度量，都可经一个邻近恒等的微分同胚化为这样的代表。
		\item 映射可以取为 \(I_0=\operatorname{Isom}(M,g_0)\) 等变；模微分同胚的光滑局部解空间，由有限维商 \(\kappa^{-1}(0)/I_0\) 描述。
	\end{enumerate}
	称 \(\kappa\) 为\term[zhangai yingshe]{障碍映射}[obstruction map]。这里 \(\kappa^{-1}(0)\) 只是有限维光滑映射的零集，不要求它为光滑流形。
\end{theorem}
\begin{proof}
	对式\eqref{g3:eq:einstein-range-equation}应用带参数隐函数定理，得到唯一 \(v=\psi(k)\)。对 \(k\) 微分，\(Lk=0\) 给出 \(\mathrm D\psi(0)=0\)。定义
	\begin{equation}\label{g3:eq:einstein-obstruction-map}
		\kappa(k)=\Pi\mathcal F(g_0+k+\psi(k)).
	\end{equation}
	由 \(\Pi L=0\) 得 \(\mathrm D\kappa(0)=0\)。原方程正好由 \(Q\) 与 \(\Pi\) 两个分量同时为零组成，因而得到第一项的集合对应。反过来，从任意小零点取 \(k=\Pi(g-g_0)\)，隐函数定理的唯一性迫使余下部分等于 \(\psi(k)\)。

	还须核对正则性与几何意义。事实上每个图上的度量，无论 \(\kappa(k)\) 是否为零，都满足 \(\mathcal F(g(k))=\kappa(k)\)。右端属于固定的光滑有限维空间。\(H^s\) 嵌入某个 \(C^{3,\alpha}\)，\(\mathcal F\) 的最高阶系数光滑依赖 \(x,g\)，所以\cref{g3:lem:elliptic-system-holder-bootstrap}给出 \(g(k)\) 光滑。对零点再用\cref{g3:thm:einstein-gauge-equivalence}，即得 Einstein 方程和散度条件。附近 Einstein 度量的代表存在性则由切片定理给出，这里输入的邻域可取 \(H^{s+1}\) 或更强的光滑邻域。

	对 \(a\in I_0\)，拉回在背景 \(L^2\) 中为等距，且与 \(L\) 交换，故保持 \(\mathcal K,X_1,Y_1\)，并与 \(\Pi,Q\) 交换。可使用由 \(g_0\) 定义的 \(I_0\) 不变 Sobolev 范数，将定义域缩为不变小球。式\eqref{g3:eq:einstein-range-equation}等变，解的唯一性遂给出 \(\psi(a^*k)=a^*\psi(k)\) 和 \(\kappa(a^*k)=a^*\kappa(k)\)。\Cref{g3:prop:metric-slice-residual-isometries}说明代表之间剩余的等价恰为 \(I_0\) 作用。

	上述对应也尊重局部拓扑。参数由连续的有限秩投影 \(\Pi\) 恢复，\(\psi\) 在 \(H^s\) 中连续。椭圆估计逐阶给出图上度量的高阶连续性：若 \(k_j\to k\)，则右端 \(\kappa(k_j)\) 在每阶都收敛；以已知低阶界对方程差使用内部估计，在稍小坐标片上逐阶升级，并由有限覆盖合并。故 \(g(k_j)\to g(k)\) 光滑。结合切片的连续性得到所述商空间的局部描述。
\end{proof}

\begin{figure}[htbp]
	\centering
	% !TeX root = ../Book3.tex
\begin{tikzcd}[row sep=2.2em,column sep=3.2em]
	(k,v)\in\mathcal K\oplus X_1
	\arrow[r,"Q\mathcal F=0"] & v=\psi(k)
	\arrow[r,"\Pi\mathcal F"] & \kappa(k)\in\mathcal K \\
	& \kappa^{-1}(0)
	\arrow[u,hook]
	\arrow[r,"\text{除以 }I_0"'] & \kappa^{-1}(0)/I_0
\end{tikzcd}

	\caption[Einstein 方程的有限维约化]{先解可逆的值域分量，得到图 \(v=\psi(k)\)；图上剩余的 \(\mathcal K\) 分量是障碍。最后的等距群商处理代表之间的剩余等价，与求解障碍方程是两个步骤。}
	\label{g3:fig:einstein-obstruction-reduction}
\end{figure}

\begin{proposition}\label{g3:prop:einstein-second-order-obstruction}
	设 \((M^n,g_0)\) 为闭连通光滑 Einstein 流形，\(n\ge3\)、\(\operatorname{Ric}_{g_0}=\lambda g_0\)，取整数 \(s>n/2+3\)。令 \(\mathcal F(g)=\operatorname{Ric}_g-\tfrac12s_gg+\tfrac{n-2}2\lambda g+\delta_g^*\delta_{g_0}(g-g_0)\)、\(L=\mathrm D\mathcal F|_{g_0}\)、\(\mathcal K=\ker L\)，\(\Pi\) 为到 \(\mathcal K\) 的 \(L^2(g_0)\) 投影、\(Q=I-\Pi\)，\(G=(L|_{\mathcal K^\perp})^{-1}\)。取\cref{g3:thm:einstein-local-obstruction}构造的核正交解映射 \(\psi\) 与障碍映射 \(\kappa\)，令 \(B=\mathrm D^2\mathcal F|_{g_0}\)。则
	\begin{equation}\label{g3:eq:einstein-quadratic-obstruction}
		\begin{aligned}
		\psi(k)&=-\tfrac12GQ B[k,k]+O(\|k\|^3),\\
		\kappa(k)&=\tfrac12\Pi B[k,k]+O(\|k\|^3).
		\end{aligned}
	\end{equation}
	若一条二次可微的同常数 Einstein 度量曲线经规范固定后在零点的速度为 \(k\in\mathcal K\)，则 \(\Pi B[k,k]=0\)。
\end{proposition}
\begin{proof}
	Taylor 展开给出 \(\mathcal F(g_0+h)=Lh+B[h,h]/2+O(\|h\|_{H^s}^3)\) 于 \(H^{s-2}\)。由 \(\psi(0)=\mathrm D\psi(0)=0\) 有 \(\psi(k)=O(\|k\|^2)\)。代入值域分量，得到
	\(L\psi(k)+QB[k,k]/2=O(\|k\|^3)\)，用 \(G\) 得第一式；用 \(\Pi L=0\) 得第二式。对题设曲线，规范选择的 \(C^1\) 依赖已保证规范后核坐标 \(k(t)=tk+o(t)\)。由 \(\kappa(k(t))=0\) 以及第二式，除以 \(t^2\) 并令 \(t\to0\)，即得 \(\Pi B[k,k]=0\)。这里无需假定规范后曲线仍二次可微。
\end{proof}

因此二阶障碍可逐个用数值配对 \(\langle B[k,k],k_a\rangle_0\) 检验。二阶项消失仍可能留下三阶或更高阶障碍；只有整个 \(\kappa\) 消失，才得到由全部 \(k\) 参数化的解族。例如有限维方程 \((x,y)\mapsto(xy,y^2)\) 在原点线性化为零，但其零集只为 \(y=0\)，说明核的维数不能直接替代实际解集的维数。

\ProseParagraphBreak

若允许 Einstein 常数变化，令 \(\rho\) 为额外的实参数，将式\eqref{g3:eq:einstein-conserved-gauge}中的 \(\lambda\) 换成 \(\lambda+\rho\)。同一个核正交逆仍可解值域分量，得到
\begin{equation}\label{g3:eq:einstein-variable-constant-model}
	\psi(k,\rho),\qquad
	\widetilde\kappa(k,\rho)=\Pi\mathcal F_{\lambda+\rho}
	\bigl(g_0+k+\psi(k,\rho)\bigr).
\end{equation}
这是 \(\mathcal K\times\mathbb R\to\mathcal K\) 的有限维方程，其零点描述所有邻近 Einstein 常数。证明逐字采用带参数隐函数定理；守恒恒等式对每个常数 \(\lambda+\rho\) 都成立，所以规范消除也仍成立。具体地，
\[
	\partial_\rho\widetilde\kappa(0,0)
	=\Pi\bigl(\tfrac{n-2}{2}g_0\bigr).
\]
若还要求体积为 \(V_0=\operatorname{Vol}_{g_0}(M)\)，就在此有限维系统中再加
\[
	\operatorname{Vol}_{g_0+k+\psi(k,\rho)}(M)=V_0.
\]
这给出固定体积、未预先固定 Einstein 常数的局部模型，仍须除以 \(I_0\)。当 \(\lambda\ne0\) 时，邻近的常数 \(\lambda'\) 同号，缩放 \((\lambda'/\lambda)g\) 将 \(\operatorname{Ric}_g=\lambda'g\) 化为常数 \(\lambda\)；当 \(\lambda=0\) 时不能作这个除法，应保留上面的参数方程。对固定 \(\lambda=0\) 的 Ricci 平坦解，常数缩放则直接提供唯一的单位体积代表。
% 纲要条目排为编号小节标题，并进入目录与书签；说明段仍保留于纲要注释中。

% !TeX root = ../../Book3.tex
% 纲要标题：### 25.4 模型和几何限制
% 纲要位置：几何学计划纲要.md，第 3190 行。
% 正文状态：初稿；尚未正式审核。

% 纲要内容开始
% 
% * 空间形式、乘积及 Ricci 平坦模型
% * 自同构导致的模空间奇点
% * Einstein 度量的退化与 Ricci 平坦存在性
% 
% ---
% 
% 纲要内容结束

\section{模型和几何限制}
\label{g3:sec:2504}

局部障碍模型以一个已经存在的 Einstein 度量为中心。具体模型可以进一步回答三个问题：核中的方向是否真正产生邻近度量，取微分同胚商时是否还有有限对称需要识别，以及一族解离开局部邻域后可能怎样退化。本节从常数缩放、乘积和环面出发，把这些现象分别算出；最后说明附加复结构怎样提供一种整体存在性结果。

\subsection{空间形式、乘积及 Ricci 平坦模型}
\label{g3:subsec:2504:01}
同一个度量的整体尺度会改变 Einstein 常数，却不改变它的 Levi--Civita 联络。先把这个区别写清楚，才能比较不同归一化下的解。

\begin{proposition}\label{g3:prop:einstein-scaling-product}
	设 \((M^n,g)\) 为光滑黎曼流形，\(n\ge3\)，且 \(\operatorname{Ric}_g=\lambda g\)。对常数 \(c>0\)，有
	\[
		\operatorname{Ric}_{cg}=\operatorname{Ric}_g
		=\frac\lambda c(cg),\qquad
		s_{cg}=c^{-1}s_g,\qquad
		\mu_{cg}=c^{n/2}\mu_g.
	\]
	再设 \((M_i^{n_i},g_i)\) 为光滑 Einstein 流形，\(\operatorname{Ric}_{g_i}=\lambda_i g_i\)、\(n_i\ge1\)，且 \(n_1+n_2\ge3\)。对 \(a,b>0\)，乘积度量 \(g=a g_1\oplus b g_2\) 为 Einstein 度量，当且仅当
	\begin{equation}\label{g3:eq:einstein-product-scales}
		\frac{\lambda_1}{a}=\frac{\lambda_2}{b}.
	\end{equation}
	此时两边的公共值就是 \(g\) 的 Einstein 常数。
\end{proposition}
\begin{proof}
	Christoffel 公式中的逆度量贡献 \(c^{-1}\)，度量的一阶导数贡献 \(c\)，故常数缩放不改变联络及其 \((1,3)\) 型曲率。Ricci 张量由这一曲率作一上一下的收缩，因而作为 \((0,2)\) 型张量也不变。标量曲率还须用逆度量缩并一次，于是乘以 \(c^{-1}\)；体积密度的变化由 \(\det(cg)=c^n\det g\) 得到。

	乘积度量的联络分别作用在两个因子上，混合方向的协变导数及曲率为零。因此
	\[
		\operatorname{Ric}_{a g_1\oplus b g_2}
		=\lambda_1g_1\oplus\lambda_2g_2.
	\]
	将 Einstein 方程分别限制到两个非零切空间，恰好得到式\eqref{g3:eq:einstein-product-scales}。
\end{proof}

若两个因子的 Einstein 常数同号且非零，可以调节相对尺度使式\eqref{g3:eq:einstein-product-scales}成立；异号时不能如此。一项为零、另一项非零也不能得到 Einstein 乘积。两个因子都 Ricci 平坦时，任意正相对尺度都给出解。这说明相对缩放有时是真实的形变参数，有时则被方程固定。

\begin{example}\label{g3:exm:einstein-sphere-product}
	设 \(p,q\ge2\)，\(g_{S^p}\)、\(g_{S^q}\) 为截面曲率一的圆球度量。给定 \(\lambda>0\)，
	\[
		g=\frac{p-1}{\lambda}g_{S^p}
		\oplus\frac{q-1}{\lambda}g_{S^q}
	\]
	满足 \(\operatorname{Ric}_g=\lambda g\)。两个因子内部的截面曲率分别为 \(\lambda/(p-1)\)、\(\lambda/(q-1)\)，混合二维平面的截面曲率为零。因此 Einstein 条件并不要求每个二维方向具有相同曲率。
\end{example}

对于截面曲率恒为 \(\kappa\) 的空间形式，式\eqref{g3:eq:constant-curvature-tensor}给出
\[
	\operatorname{Ric}=(n-1)\kappa g,
	\qquad
	\mathring R(h)=\kappa\bigl((\operatorname{tr}_g h)g-h\bigr).
\]
第二式是把该曲率张量与对称二张量 \(h\) 缩并。它让正曲率空间形式的局部问题化为一个带正零阶项的能量恒等式。

\begin{corollary}\label{g3:cor:positive-space-form-einstein-rigid}
	设 \((M^n,g_0)\) 为闭连通光滑黎曼流形，\(n\ge3\)，截面曲率恒为 \(\kappa>0\)。固定 \(\lambda=(n-1)\kappa\)。则 \(g_0\) 在 Einstein 常数为 \(\lambda\) 的度量中模微分同胚局部刚性。
\end{corollary}
\begin{proof}
	由\cref{g3:prop:einstein-kernel-tt}，规范线性化的核由满足 \(\delta_{g_0}h=0\)、\(\operatorname{tr}_{g_0}h=0\) 的张量组成，且每个核元素满足
	\[
		0=(\nabla^*\nabla-2\mathring R)h
		=\nabla^*\nabla h+2\kappa h.
	\]
	与 \(h\) 作 \(L^2\) 内积得到 \(\|\nabla h\|_2^2+2\kappa\|h\|_2^2=0\)，故 \(h=0\)。应用\cref{g3:cor:einstein-infinitesimal-rigidity}便得结论。
\end{proof}

这里固定 Einstein 常数消除了整体缩放。若 \(\kappa<0\)，同一算式中的零阶项改变符号，刚性需要另作论证，不能由上面的非负积分直接读出。

\ProseParagraphBreak

Ricci 平坦环面给出相反的情形：核的每个方向都能积分成一族精确解。取 \(T^n=\mathbb R^n/\mathbb Z^n\)，\(n\ge3\)，对每个正定实对称矩阵 \(Q\) 定义
\begin{equation}\label{g3:eq:flat-torus-matrix-metric}
	g_Q=\sum_{i,j=1}^nQ_{ij}\,\mathrm dx^i\otimes\mathrm dx^j,
	\qquad \mathcal P_n=\{Q=Q^{\mathsf T}>0\}.
\end{equation}
系数为常数，故联络系数及曲率都为零；基本立方体的体积为 \(\sqrt{\det Q}\)。因此单位体积子族为
\[
	\mathcal P_n^1=\{Q\in\mathcal P_n:\det Q=1\},
	\qquad \dim\mathcal P_n^1=\frac{n(n+1)}2-1.
\]
行列式的一阶变分为 \(D(\det)_QH=\det Q\operatorname{tr}(Q^{-1}H)\)，所以最后一个维数公式也说明固定体积如何去掉一个缩放方向。

\begin{proposition}\label{g3:prop:flat-torus-einstein-kernel}
	记 \(T^n=\mathbb R^n/\mathbb Z^n\)，\(\mathcal P_n\) 为正定实对称矩阵空间，\(g_Q=\sum Q_{ij}\,\mathrm dx^i\mathrm dx^j\)，\(\mathcal P_n^1=\{Q\in\mathcal P_n:\det Q=1\}\)。
	设 \(n\ge3\)，\(Q\in\mathcal P_n\)，在闭环面 \((T^n,g_Q)\) 上取 Einstein 常数 \(\lambda=0\) 及本章的散度规范。规范线性化的核恰为常系数对称二张量，维数为 \(n(n+1)/2\)。每个核方向均由常系数 Ricci 平坦度量族实现；局部障碍映射在这一模型中恒为零。固定体积后的切空间再满足 \(\operatorname{tr}(Q^{-1}H)=0\)。
\end{proposition}
\begin{proof}
	由\cref{g3:prop:einstein-kernel-tt}，核元素的迹为常数，去掉其常数倍 \(g_Q\) 后得到满足 \(\nabla^*\nabla h=0\) 的无迹张量。闭流形上的分部积分给出 \(\nabla h=0\)。平坦坐标中的平行张量正是常系数张量。反过来，常系数张量的所有导数为零，而曲率及 \(\lambda\) 也为零，代入线性化便知它属于核。

	任意对称常矩阵 \(H\) 在充分小的 \(|t|\) 下给出正定矩阵 \(Q+tH\)，相应度量都是 Ricci 平坦的，并满足相对于 \(g_Q\) 的散度规范。更一般地，充分小的核元素本身就给出一个精确解；在\cref{g3:thm:einstein-local-obstruction}的核与正交补分解中，这些解的正交补分量为零。该分量的隐函数唯一性遂说明约化图的正交补函数和障碍映射都恒为零。

	若 \(\det Q=1\) 且 \(\operatorname{tr}(Q^{-1}H)=0\)，可用
	\[
		Q(t)=Q^{1/2}\exp\bigl(tQ^{-1/2}HQ^{-1/2}\bigr)Q^{1/2}
	\]
	实现给定切方向。矩阵指数保持正定性，\(Q'(0)=H\)，且 \(\det Q(t)=\det Q\exp(t\operatorname{tr}(Q^{-1}H))=1\)。
\end{proof}

矩阵 \(Q\) 记录的是选定整数坐标后的度量。要得到几何对象的参数，还必须识别改换整数坐标所产生的相同环面。

\begin{proposition}\label{g3:prop:flat-torus-matrix-moduli}
	记 \(T^n=\mathbb R^n/\mathbb Z^n\)，\(\mathcal P_n\) 为正定实对称矩阵空间，\(g_Q=\sum Q_{ij}\,\mathrm dx^i\mathrm dx^j\)，\(\mathcal P_n^1=\{Q\in\mathcal P_n:\det Q=1\}\)。
	设 \(n\ge3\)、\(Q,Q'\in\mathcal P_n\)。常系数平坦环面 \((T^n,g_Q)\) 与 \((T^n,g_{Q'})\) 等距，当且仅当存在 \(A\in\operatorname{GL}(n,\mathbb Z)\) 使
	\[
		Q'=A^{\mathsf T}QA.
	\]
	因此单位体积的常系数平坦环面等距类由 \(\mathcal P_n^1\) 在上述整数矩阵作用下的轨道表示。
\end{proposition}
\begin{proof}
	若有这样的 \(A\)，映射 \([x]\mapsto[Ax]\) 为环面微分同胚，其拉回度量为 \(g_{Q'}\)。反过来，一个环面等距映射可以提升为两个万有覆盖间的等距映射 \(F:\mathbb R^n\to\mathbb R^n\)。等距映射保持 Levi--Civita 联络，而两边的坐标联络系数都为零，因此 \(\partial_i\partial_jF=0\)，即 \(F(x)=Ax+b\)。它与覆盖变换相容，故 \(A\mathbb Z^n=\mathbb Z^n\)；其逆也满足同一条件，所以 \(A\in\operatorname{GL}(n,\mathbb Z)\)。等距等式正是 \(Q'=A^{\mathsf T}QA\)。平移 \(b\) 不改变常系数度量。
\end{proof}

\subsection{自同构导致的模空间奇点}
\label{g3:subsec:2504:02}
前一命题中的平移已经显示：一个度量可以有连续的等距群，而该群在某些形变参数上作用平凡。判断商是否出现奇性，需要检查稳定子在参数空间上的实际表示。

先固定 \(Q\in\mathcal P_n^1\)。其整数稳定子为
\[
	\Gamma_Q=\{A\in\operatorname{GL}(n,\mathbb Z):A^{\mathsf T}QA=Q\}.
\]
它是有限群。事实上，若 \(mI\le Q\le MI\)，则其任一元素的第 \(i\) 列满足 \(m|Ae_i|^2\le Q_{ii}\le M\)，整数列向量只有有限种。类似地，对于 \(Q,Q'\) 同时落在 \(\mathcal P_n^1\) 的一个紧集内，只可能有有限多个整数矩阵满足 \(Q'=A^{\mathsf T}QA\)。这证明作用在紧集上的有限性。缩小 \(Q\) 的邻域，可以排除不属于 \(\Gamma_Q\) 的这些矩阵，局部商于是由有限稳定子决定。

令
\[
	V_Q=\{H=H^{\mathsf T}:\operatorname{tr}(Q^{-1}H)=0\}.
\]
稳定子在此空间上的作用为 \(H\mapsto A^{\mathsf T}HA\)。这个线性表示给出实际的局部商坐标：对充分小的 \(H\in V_Q\)，令
\[
	\chi(H)=\frac{Q+H}{\det(Q+H)^{1/n}}.
\]
有 \(\chi(0)=Q\)、\(D\chi_0=I\)，故有限维逆函数定理使 \(\chi\) 成为局部坐标。行列式在整数合同变换下不变，因而 \(\chi\) 对 \(\Gamma_Q\) 等变。局部等距类空间就是这一坐标球的有限群商。特别地，\(-I\) 对全部对称矩阵作用恒等；它作为等距映射仍是非平凡的，但在这里属于表示的核。

\begin{example}\label{g3:exm:flat-torus-stabilizer-cone}
	在 \(n=3\) 时，选实数 \(a,b,c\)，使 \(b>a^2\)、\(c>0\)，且 \(1,a,b,c\) 在 \(\mathbb Q\) 上线性无关。这样的选择存在于相应开集内，因为所有非零有理线性关系只给出可数个真超平面。令
	\[
		Q_0=\begin{pmatrix}1&a&0\\a&b&0\\0&0&c\end{pmatrix},
		\qquad Q=(\det Q_0)^{-1/3}Q_0.
	\]
	若整数矩阵 \(A\) 保持 \(Q_0\)，将其三行分别记为 \(r_1,r_2,r_3\)，则
	\[
		A^{\mathsf T}Q_0A
		=r_1^{\mathsf T}r_1
		+a(r_1^{\mathsf T}r_2+r_2^{\mathsf T}r_1)
		+b r_2^{\mathsf T}r_2+c r_3^{\mathsf T}r_3.
	\]
	逐项比较 \(1,b,c\) 的有理系数可知 \(r_i=\varepsilon_i e_i^{\mathsf T}\)，\(\varepsilon_i\in\{1,-1\}\)；再比较 \(a\) 的系数得 \(\varepsilon_1\varepsilon_2=1\)。所以
	\[
		\Gamma_Q=\{\operatorname{diag}(\varepsilon,\varepsilon,\eta):
		\varepsilon,\eta\in\{1,-1\}\}.
	\]
	除去作用平凡的 \(\{I,-I\}\) 后，有效稳定子由 \(D=\operatorname{diag}(-1,-1,1)\) 生成。把切向量写成
	\[
		H=\begin{pmatrix}B&v\\v^{\mathsf T}&t\end{pmatrix},
		\qquad B=B^{\mathsf T}\in\mathbb R^{2\times2},\quad v\in\mathbb R^2,
	\]
	则 \(D^{\mathsf T}HD\) 保持 \(B,t\)，把 \(v\) 变为 \(-v\)。迹约束只限制 \(B,t\)，故不变方向有三维，变号方向有二维。单位体积模族在 \([Q]\) 附近具有局部模型
	\begin{equation}\label{g3:eq:flat-torus-local-cone}
		\mathbb R^3\times\bigl(\mathbb R^2/\{v\sim-v\}\bigr).
	\end{equation}
	给切空间取不变内积，二维变号因子的商度量是总角为 \(\pi\) 的平锥。在原点的小圆周上，周长与半径之比为 \(\pi\)，而非 Euclid 平面的 \(2\pi\)。因此稳定子在自然局部商结构中留下锥角。其底层拓扑空间仍同胚于 \(\mathbb R^5\)：二维商的同胚可由 \(z\mapsto z^2\) 给出。保留稳定子与商度量，比只问底层集合是否像一个开球记录了更多几何信息。
\end{example}

在一般 Einstein 局部模型中还可能出现另一种奇性：障碍映射的零集本身就不光滑。\Cref{g3:thm:einstein-local-obstruction}先得到有限维零集，再识别剩余等距群作用。核维数、零集的局部形状和稳定子表示分别属于三个层次；平坦环面中障碍消失，式\eqref{g3:eq:flat-torus-local-cone}的现象完全来自最后一层。

\subsection{Einstein 度量的退化与 Ricci 平坦存在性}
\label{g3:subsec:2504:03}
有限维局部模型只描述某个解的邻域。即使方程和体积归一化都固定，解族仍可能在距离尺度上逃出所有这样的邻域。

\begin{example}\label{g3:exm:unit-volume-flat-torus-degeneration}
	设 \(n\ge3\)、\(0<\varepsilon<1\)，在 \(T^n=\mathbb R^n/\mathbb Z^n\) 上取
	\[
		g_\varepsilon=\varepsilon^2(\mathrm dx^1)^2
		+\varepsilon^{-2}(\mathrm dx^2)^2
		+\sum_{i=3}^n(\mathrm dx^i)^2.
	\]
	这是单位体积的平坦 Einstein 度量，且
	\[
		\operatorname{inj}(g_\varepsilon)=\frac\varepsilon2,
		\qquad
		\operatorname{diam}(g_\varepsilon)
		=\frac12\sqrt{\varepsilon^2+\varepsilon^{-2}+n-2}.
	\]
	第一个等式来自最短非零格向量的长度为 \(\varepsilon\)；没有共轭点时，注入半径是这一长度的一半。第二个等式来自乘积圆周的距离平方为各因子距离平方之和，每个因子的最大距离恰为其周长的一半。于是曲率恒为零、体积恒为一，直径却趋于无穷。任意拉回都保持直径，因此这列度量不能在模微分同胚后光滑收敛到一个闭流形上的黎曼度量。它也不能有紧度量空间的 GH 收敛子列：若对应的畸变趋零，则两空间的直径之差也趋零。
\end{example}

若保持其余圆周长度为一，只压缩第一个方向，则得到
\[
	\widetilde g_\varepsilon
	=\varepsilon^2(\mathrm dx^1)^2+\sum_{i=2}^n(\mathrm dx^i)^2.
\]
此时直径有界、体积趋零。到单位平坦 \(T^{n-1}\) 的投影是满射，任意一对点的距离损失不超过 \(\varepsilon/2\)，故由\cref{g3:thm:gh-correspondence}有
\[
	d_{\mathrm{GH}}\bigl((T^n,\widetilde g_\varepsilon),T^{n-1}\bigr)
	\le\frac\varepsilon4.
\]
这是\cref{g3:prop:flat-torus-collapse}的高维乘积版本。前一族用长方向补偿短方向，后一族允许总体积消失；体积归一化与直径控制在紧性问题中承担不同作用。

\ProseParagraphBreak

整体存在性还可能来自额外的几何结构。这里给出一个精确的外部结果，说明 Ricci 平坦度量的来源远不止常系数环面。回顾\cref{g3:def:parallel-kahler-structure}，Kähler 结构使正交复结构 \(J\) 平行，基本二形式 \(\omega(X,Y)=g(JX,Y)\) 闭。在复流形表述中，固定复结构后，正的实闭 \((1,1)\) 形式给出这样的度量；\([\omega]\in H^2(X;\mathbb R)\) 称为其 Kähler 类。

复切丛的\emph{实第一陈类}是第一陈类在 \(H^2(X;\mathbb R)\) 中的像。对于 Kähler 度量，它由 Ricci 形式的 \(1/(2\pi)\) 倍表示。在全纯坐标中，若 \(\omega=\mathrm i\sum g_{\alpha\bar\beta}\,\mathrm dz^\alpha\wedge\mathrm d\bar z^\beta\)，则 Ricci 形式为
\[
	\rho_g=-\mathrm i\,\partial\bar\partial\log\det(g_{\alpha\bar\beta}),
	\qquad \rho_g(X,Y)=\operatorname{Ric}_g(JX,Y).
\]
这里 \(\partial\)、\(\bar\partial\) 分别取外微分中提高全纯次数、反全纯次数的部分。更换同一复结构上的 Kähler 度量时，两式之差是对整体体积密度比的对数作用 \(-\mathrm i\partial\bar\partial\)，因而为恰当形式。这解释了上述上同调类为何不随度量改变；它也说明 Ricci 平坦时实第一陈类必须为零。

\begin{theorem}\label{g3:thm:kahler-ricci-flat-existence-input}
	设 \(X\) 为紧连通复流形，复维数 \(m\ge2\)，并设 \(X\) 有 Kähler 度量且实第一陈类 \(c_1(X)_{\mathbb R}=0\)。则每个 Kähler 类中恰有一个 Ricci 平坦的 Kähler 度量。
\end{theorem}
该定理作为外部输入，采用\cite[定理1，第1798页]{Yau1977Calabi}中固定 Kähler 类、指定 Ricci 形式的结论，并把所指定形式取为零。所引文章是该存在性结果的原始公告；其全局分析依赖复 Monge--Ampère 方程的先验估计，不能由已知解附近的隐函数定理单独推出。唯一性固定了复结构和 Kähler 类：改变这些数据，仍可产生不同的 Ricci 平坦度量。文献\cite[定理2，第1798页]{Yau1977Calabi}还给出紧单连通、Ricci 平坦而完整曲率不恒为零的 Kähler 流形，因而 Ricci 平坦也不等同于平坦。

这些模型把本章局部结论放在了具体的位置：正曲率空间形式的核消失，平坦环面的核全部可积，有限稳定子使光滑参数仍需取商，退化族则说明局部参数不能直接组成一个紧的整体分类空间。
% 纲要条目排为编号小节标题，并进入目录与书签；说明段仍保留于纲要注释中。


\begin{chaptersummary}
	度量的切向量是对称二张量。Ricci 线性化由 Lichnerowicz 型算子与 \(-\delta^*\beta\) 组成；后者反映微分同胚方向的退化。标量曲率变分再经分部积分，给出 Einstein--Hilbert 泛函在固定体积下的临界方程。闭性、维数 \(n\ge3\) 以及 Einstein 常数的固定方式，在这些论证中各有明确作用。

	\ProseParagraphBreak
	守恒张量 \(\mathcal G_\lambda\) 加上 \(\delta_g^*\delta_0(g-g_0)\) 后得到椭圆系统。它的符号可逆，但不正定。对方程取散度并配对，先使规范项消失，再利用 Killing 场正交补上的估计恢复散度条件。局部切片由 \(2\delta_0\delta_0^*\) 求逆构造，复合中的一阶微分损失通过 Sobolev 空间的选择处理，光滑性由椭圆正则性恢复。

	\ProseParagraphBreak
	线性化 \(L\) 形式自伴，故其核 \(\mathcal K\) 同时表示无穷小形变与余核相容条件。隐函数定理给出 \(v=\psi(k)\)，局部解对应于有限维障碍方程 \(\kappa(k)=0\)。二阶障碍为 \(\tfrac12\Pi\mathrm D^2\mathcal F[k,k]\)；它消失后还可能出现高阶障碍。模微分同胚的局部空间是这个零点集再模背景等距群。令 Einstein 常数一同变化，只增加一个有限维参数；体积归一化则增加一条有限维约束。

	\ProseParagraphBreak
	正截面曲率空间形式提供局部刚性的例子，Ricci 平坦乘积和平坦环面提供真实形变。环面的整数线性自同构说明群作用能够产生有限商模型，而体积固定的细长环面说明局部模理论并不提供整体紧性。Kähler 类中的 Ricci 平坦存在性由 Yau 定理给出；其唯一性固定复结构和 Kähler 类，与不同类之间的形变并不冲突。
\end{chaptersummary}

\BookExercises

\begin{exercise}\label{g3:ex:einstein-volume-second-variation}
	设 \((M,g)\) 闭，\(h,k\) 是光滑对称二张量。证明总体积的二阶微分为
	\[
		\mathrm D^2\operatorname{Vol}_g(h,k)
		=\int_M\left(\tfrac14\operatorname{tr}_gh\operatorname{tr}_gk
		-\tfrac12\langle h,k\rangle_g\right)\mu_g.
	\]
	若 \(\operatorname{tr}_gh=0\) 且 \(h\ne0\)，说明直线 \(g+th\) 为什么不是固定体积的曲线，并求曲线 \(g_t=g+th+t^2a/2+o(t^2)\) 保持体积到二阶所需的积分条件。
\end{exercise}
\begin{solution}
	由逆矩阵微分 \(\mathrm D(g^{-1})k=-g^{-1}kg^{-1}\)，有 \(\mathrm D(\operatorname{tr}_gh)k=-\langle h,k\rangle_g\)。再对 \(\mathrm D\operatorname{Vol}_g(h)=\tfrac12\int\operatorname{tr}_gh\,\mu_g\) 微分，用 \(\mathrm D\mu_g(k)=\tfrac12\operatorname{tr}_gk\,\mu_g\)，便得所列公式。无迹时直线的体积一阶导数为零，二阶导数为 \(-\tfrac12\|h\|_2^2<0\)。一般曲线还贡献 \(\mathrm D\operatorname{Vol}_g(a)\)，故二阶体积不变要求
	\[
		\int_M\operatorname{tr}_ga\,\mu_g=\int_M|h|_g^2\,\mu_g.
	\]
	切空间中的体积约束只是一阶条件；构造约束曲线时必须补上二阶修正。
\end{solution}

\begin{exercise}\label{g3:ex:einstein-sphere-conformal-gauge}
	在截面曲率为 \(\kappa>0\) 的圆球 \((S^n,g)\)、\(n\ge3\) 上，令 \(\lambda=(n-1)\kappa\)，\(f\) 为非零高度函数。证明 \(h=fg\) 满足 \(\mathrm D(\operatorname{Ric}-\lambda g)_g h=0\)，但它是纯微分同胚方向，并且不满足背景散度条件。
\end{exercise}
\begin{solution}
	将 \(h=fg\) 代入曲率线性化，得
	\[
		\mathrm D(\operatorname{Ric}-\lambda g)_g(fg)
		=-\tfrac{n-2}{2}\operatorname{Hess}f
		+\left(\tfrac12\Delta f-\lambda f\right)g.
	\]
	高度函数满足 \(\operatorname{Hess}f=-\kappa fg\)、\(\Delta f=n\kappa f\)，右端系数相加为零。取 \(X=-(2\kappa)^{-1}\operatorname{grad}f\)，则 \(\mathcal L_Xg=2\operatorname{Hess}(-(2\kappa)^{-1}f)=fg\)。另一方面 \(\delta(fg)=-\mathrm df\ne0\)。因此原方程的非零线性化核不必代表真正的模参数，这与圆球的局部刚性相容。
\end{solution}

\begin{exercise}\label{g3:ex:einstein-pure-gauge-energy}
	设 \((M,g_0)\) 闭且 Einstein，\(L\) 为式\eqref{g3:eq:einstein-gauged-linearization}的算子。对 \(h=\mathcal L_Xg_0\)，证明
	\[
		Lh=\delta_0^*\delta_0h,\qquad
		\langle Lh,h\rangle_0=\|\delta_0h\|_2^2.
	\]
	证明 \(Lh=0\) 当且仅当 \(X\) 为 Killing 场。
\end{exercise}
\begin{solution}
	由 \(\mathcal G_\lambda\) 的自然性以及 \(\mathcal G_\lambda(g_0)=0\)，沿 \(X\) 的流微分得 \(\mathcal G_\lambda'(\mathcal L_Xg_0)=\mathcal L_X\mathcal G_\lambda(g_0)=0\)。代入 \(L\) 并分部积分，得到两式。若 \(Lh=0\)，则 \(\delta_0h=0\)；又 \(h=2\delta_0^*X^\flat\)，故 \(\|h\|_2^2=2\langle\delta_0h,X^\flat\rangle_0=0\)，从而 \(X\) 为 Killing 场。反之 Killing 场给 \(h=0\)。这里的非负性只针对纯规范方向，不能推得 \(L\) 在全部张量上非负。
\end{solution}

\begin{exercise}\label{g3:ex:einstein-symbol-inverse}
	设 \((M^n,g_0)\) 为光滑 Einstein 背景，\(n\ge3\)，\(L\) 为式\eqref{g3:eq:einstein-gauged-linearization}的规范线性化，\(h,q\) 为 \(x\in M\) 处的对称二张量。对单位余向量 \(\xi=e^1\)，设 \(q=\sigma_L(\xi)h\)。由式\eqref{g3:eq:einstein-gauge-block-symbol}显式求 \(h\)，并解释维数二为何例外。
\end{exercise}
\begin{solution}
	令 \(a,b\ge2\)、\(H=\sum_{a=2}^nh_{aa}\)，则
	\[
		h_{11}=q_{11},\qquad h_{1a}=2q_{1a},\qquad
		H=-\frac{2}{n-2}\sum_{a=2}^nq_{aa},\qquad
		h_{ab}=2q_{ab}+H\delta_{ab}.
	\]
	这些式子给出唯一逆像。维数二时横向只有一个分量，\(h_{22}-H=0\) 恒成立；例如仅有 \(h_{22}=1\) 的张量被符号消去。因此这一规范化系统在二维不椭圆。几何上二维 \(\operatorname{Ric}-\tfrac12sg\) 恒为零，所用守恒张量已经失去原来的方程信息。
\end{solution}

\begin{exercise}\label{g3:ex:einstein-ricci-flat-variable-constant}
	设 \(g_t\) 是闭 \(n\) 维流形上可微的光滑度量族，满足 \(\operatorname{Ric}_{g_t}=\lambda(t)g_t\)，且 \(g_0\) Ricci 平坦。证明 \(\lambda'(0)=0\)。
\end{exercise}
\begin{solution}
	设 \(h=\dot g_0\)。从 \(s_{g_t}=n\lambda(t)\) 得 \(s'_0=n\lambda'(0)\)。在 Ricci 平坦背景下，标量曲率变分为 \(s'_0=\Delta_0\operatorname{tr}_0h+\delta_0\delta_0h\)，两项在闭流形上的积分均为零。因此 \(n\lambda'(0)\operatorname{Vol}_{g_0}(M)=0\)，所求成立。这只限制曲线的一阶变化，不能据此断言附近每个 Einstein 度量的常数均为零。
\end{solution}

\begin{exercise}\label{g3:ex:einstein-ricci-flat-product-shape}
	设 \((M_1^p,g_1)\)、\((M_2^q,g_2)\) 为闭 Ricci 平坦流形，\(p,q\ge1\)、\(p+q\ge3\)。证明
	\[
		g_t=e^{2qt}g_1\oplus e^{-2pt}g_2
	\]
	是一族等体积 Ricci 平坦度量，其初速度为非零的无迹无散度张量，且不是纯微分同胚方向。
\end{exercise}
\begin{solution}
	常数缩放不改变 Ricci 张量作为二协变张量的表达，乘积 Ricci 张量按因子分解，所以每个 \(g_t\) 都 Ricci 平坦。两个体积因子分别为 \(e^{pqt}\)、\(e^{-pqt}\)，相乘为一。初速度 \(h=2qg_1\oplus(-2pg_2)\) 平行，因而散度为零；其迹为 \(2pq-2pq=0\)，且 \(h\ne0\)。若 \(h=\mathcal L_Xg_0\)，则 \(\|h\|_2^2=2\langle\delta_0h,X^\flat\rangle_0=0\)，矛盾。因此相对尺度是真正的无穷小形变，并由所给曲线积分。
\end{solution}

\begin{exercise}\label{g3:ex:einstein-flat-torus-reparametrization}
	在 \(T^n=\mathbb R^n/\mathbb Z^n\)、\(n\ge3\) 上，令
	\[
		A_m=\begin{pmatrix}1&m\\0&1\end{pmatrix}\oplus I_{n-2},\qquad
		Q_m=A_m^{\mathsf T}A_m,\qquad m\in\mathbb Z.
	\]
	证明矩阵 \(Q_m\) 的系数无界，但对应的所有平坦度量彼此等距。求它们的体积、注入半径和直径，并与\cref{g3:exm:unit-volume-flat-torus-degeneration}比较。
\end{exercise}
\begin{solution}
	二维块为 \(\left(\begin{smallmatrix}1&m\\m&m^2+1\end{smallmatrix}\right)\)，故系数无界。但 \(A_m\in GL(n,\mathbb Z)\)，映射 \([x]\mapsto[A_mx]\) 是环面微分同胚，且将标准度量拉回为 \(g_{Q_m}\)。因此体积恒为一，注入半径恒为 \(1/2\)，直径恒为 \(\sqrt n/2\)。这一序列在模空间中是一个点；细长环面的注入半径趋于零、直径趋于无穷，才是不能由重参数化消除的几何退化。
\end{solution}

\begin{exercise}\label{g3:ex:einstein-cubic-obstruction-model}
	对有限维映射 \(F:\mathbb R^2\to\mathbb R^2\)、\(F(x,y)=(y+x^2,xy)\)，在原点实行核与余核约化。求障碍函数，说明二阶障碍消失为何不足以保证核方向可积。
\end{exercise}
\begin{solution}
	\(\mathrm DF(0)(a,b)=(b,0)\)，故核是 \(x\) 轴，像是目标第一坐标轴，余核取第二坐标轴。像方向方程 \(y+x^2=0\) 给 \(y=\psi(x)=-x^2\)，代入另一分量得 \(\kappa(x)=-x^3\)。它的二阶项为零，但零点只有 \(x=0\)，因此 \(F^{-1}(0)\) 在原点附近也是单点。非零核向量不能成为非平凡解曲线的初速度；障碍首次出现在三阶。
\end{solution}

\begin{exercise}\label{g3:ex:einstein-sign-quotient}
	考虑 \(\mathbb R^2\) 上的等价关系 \((x,y)\sim(-x,-y)\)。用不变量 \(a=x^2\)、\(b=xy\)、\(c=y^2\) 描述商空间，证明它在拓扑上仍同胚于平面，并计算继承 Euclid 度量的商锥角。
\end{exercise}
\begin{solution}
	不变量满足 \(a,c\ge0\)、\(ac=b^2\)。反过来，若 \(a>0\)，取 \(x=\sqrt a\)、\(y=b/\sqrt a\)；若 \(a=0\)，则 \(b=0\)，取 \(x=0\)、\(y=\sqrt c\)。每个像恰对应一条符号轨道。此连续映射是适当的，因为 \(a+c=x^2+y^2\)，故诱导商空间到所述闭半代数集合的同胚。写极坐标 \((r,\theta\bmod\pi)\)，映射到平面极坐标 \((r,2\theta\bmod2\pi)\) 给出另一同胚，并在原点连续。半径 \(r\) 的商圆长度为 \(\pi r\)，所以锥角为 \(\pi\)，不同于 Euclid 平面的 \(2\pi\)。底层拓扑为平面、商度量呈锥形以及群作用具有非平凡稳定子，是三个不同的信息。
\end{solution}

\begin{exercise}\label{g3:ex:einstein-kahler-class-variation}
	在固定复环面 \(\mathbb C^2/(\mathbb Z^2+i\mathbb Z^2)\) 上，取周期为一的实坐标 \(x_j,y_j\)。考虑
	\[
		\omega_t=e^t\,\mathrm dx_1\wedge\mathrm dy_1
		+e^{-t}\,\mathrm dx_2\wedge\mathrm dy_2.
	\]
	证明它们给出体积相同的 Ricci 平坦 Kähler 度量，并说明这个族与固定 Kähler 类中的唯一性并不矛盾。
\end{exercise}
\begin{solution}
	相应度量为 \(e^t(\mathrm dx_1^2+\mathrm dy_1^2)+e^{-t}(\mathrm dx_2^2+\mathrm dy_2^2)\)，系数常数且正定，故平坦。形式 \(\omega_t\) 闭并与固定复结构相容，因而是 Kähler 形式。体积形式 \(\omega_t^2/2\) 恒为 \(\mathrm dx_1\wedge\mathrm dy_1\wedge\mathrm dx_2\wedge\mathrm dy_2\)。但 \(\omega_t\) 在两个坐标复环面上的积分分别为 \(e^t\)、\(e^{-t}\)，所以其上同调类随 \(t\) 改变。Yau 定理在每个固定 Kähler 类内给唯一性，不把这些不同类的度量等同。
\end{solution}

\begin{exercise}\label{g3:ex:einstein-flat-fourier-gauge}
	在标准平坦环面 \(T^n\)、\(n\ge3\) 上，设光滑对称张量 \(h\) 满足 \(\operatorname{Ric}'h=0\)。用 Fourier 级数直接证明 \(h=H+\mathcal L_Xg\)，其中 \(H\) 为常系数对称张量，\(X\) 为光滑周期向量场。要求给出 \(X\) 的非零频率系数。
\end{exercise}
\begin{solution}
	对非零频率 \(\xi=2\pi m\)，记 \(\widehat h=\widehat h(m)\)。线性化方程为
	\[
		|\xi|^2\widehat h+\xi\otimes\xi\operatorname{tr}\widehat h
		-\xi\otimes\widehat h(\xi,\cdot)
		-\widehat h(\xi,\cdot)\otimes\xi=0.
	\]
	取 \(\xi/|\xi|=e_1\) 后，横向分量给 \(\widehat h_{ab}=0\)（\(a,b>1\)），其余只含 \(\widehat h_{11}\)、\(\widehat h_{1a}\)。定义
	\[
		\widehat X(m)=\frac{1}{i|\xi|^2}
		\left(\widehat h\,\xi-
		\frac{\xi\,\widehat h(\xi,\xi)}{2|\xi|^2}\right),
		\qquad \widehat X(0)=0.
	\]
	在上述正交标架中直接计算，得 \(\widehat h=i(\xi\otimes\widehat X+\widehat X\otimes\xi)\)，也就是 Lie 导数的该频率系数。\(h\) 光滑使其 Fourier 系数任意次快速衰减，所给公式只再除以一次 \(|\xi|\)，故 \(X\) 光滑；实性由共轭对称性保证。取 \(H=\widehat h(0)\)，各频率相加得到结论。常数模式正是本章平坦环面模型中的真实线性参数。
\end{solution}
